%% Generated by Sphinx.
\def\sphinxdocclass{report}
\documentclass[letterpaper,10pt,english]{sphinxmanual}
\ifdefined\pdfpxdimen
   \let\sphinxpxdimen\pdfpxdimen\else\newdimen\sphinxpxdimen
\fi \sphinxpxdimen=.75bp\relax
\ifdefined\pdfimageresolution
    \pdfimageresolution= \numexpr \dimexpr1in\relax/\sphinxpxdimen\relax
\fi
\newdimen\sphinxremdimen\sphinxremdimen = 10pt
%% let collapsible pdf bookmarks panel have high depth per default
\PassOptionsToPackage{bookmarksdepth=5}{hyperref}

\PassOptionsToPackage{booktabs}{sphinx}
\PassOptionsToPackage{colorrows}{sphinx}

\PassOptionsToPackage{warn}{textcomp}
\usepackage[utf8]{inputenc}
\ifdefined\DeclareUnicodeCharacter
% support both utf8 and utf8x syntaxes
  \ifdefined\DeclareUnicodeCharacterAsOptional
    \def\sphinxDUC#1{\DeclareUnicodeCharacter{"#1}}
  \else
    \let\sphinxDUC\DeclareUnicodeCharacter
  \fi
  \sphinxDUC{00A0}{\nobreakspace}
  \sphinxDUC{2500}{\sphinxunichar{2500}}
  \sphinxDUC{2502}{\sphinxunichar{2502}}
  \sphinxDUC{2514}{\sphinxunichar{2514}}
  \sphinxDUC{251C}{\sphinxunichar{251C}}
  \sphinxDUC{2572}{\textbackslash}
\fi
\usepackage{cmap}
\usepackage[T1]{fontenc}
\usepackage{amsmath,amssymb,amstext}
\usepackage{babel}



\usepackage{tgtermes}
\usepackage{tgheros}
\renewcommand{\ttdefault}{txtt}



\usepackage[Bjarne]{fncychap}
\usepackage{sphinx}

\fvset{fontsize=auto}
\usepackage{geometry}


% Include hyperref last.
\usepackage{hyperref}
% Fix anchor placement for figures with captions.
\usepackage{hypcap}% it must be loaded after hyperref.
% Set up styles of URL: it should be placed after hyperref.
\urlstyle{same}

\addto\captionsenglish{\renewcommand{\contentsname}{Documentation}}

\usepackage{sphinxmessages}
\setcounter{tocdepth}{0}



\title{hrHSA}
\date{Jul 29, 2026}
\release{0.1.0}
\author{Paul Kasko, Lubov Kabo, and Ralph Kühn}
\newcommand{\sphinxlogo}{\vbox{}}
\renewcommand{\releasename}{Release}
\makeindex
\begin{document}

\ifdefined\shorthandoff
  \ifnum\catcode`\=\string=\active\shorthandoff{=}\fi
  \ifnum\catcode`\"=\active\shorthandoff{"}\fi
\fi

\pagestyle{empty}
\sphinxmaketitle
\pagestyle{plain}
\sphinxtableofcontents
\pagestyle{normal}
\phantomsection\label{\detokenize{index::doc}}


\sphinxAtStartPar
\sphinxstylestrong{hrHSA} \sphinxstyleemphasis{(high\sphinxhyphen{}resolution Habitat Selection Analysis)} is a scalable Python framework for rigorous statistical inference on spatially distributed observations. It integrates positional data with environmental covariates to quantify how heterogeneous spatial fields, individual variation, and changing conditions shape the probability distribution of observed locations. The framework combines resource\sphinxhyphen{}selection analysis, spatial point\sphinxhyphen{}process models, hierarchical Bayesian inference, and mechanistic simulation, with explicit treatment of availability, observation processes, uncertainty, validation, and prediction.

\sphinxAtStartPar
Developed primarily for wildlife telemetry and spatial ecology, \sphinxstylestrong{hrHSA} addresses the more general problem of inferring interactions between actively and/or passively redistributing entities and their spatial context. Its methods therefore extend naturally to applications in geospatial science, epidemiology, environmental modelling, human mobility, and other domains concerned with spatially structured processes. Reproducible, vectorized, out\sphinxhyphen{}of\sphinxhyphen{}core, and HPC\sphinxhyphen{}ready workflows support analyses from covariate extraction and parameter estimation to predictive simulation.

\sphinxstepscope


\chapter{Getting started}
\label{\detokenize{getting_started:getting-started}}\label{\detokenize{getting_started::doc}}
\sphinxAtStartPar
This guide introduces the core \sphinxstylestrong{frequentist habitat\sphinxhyphen{}selection workflow} in hrHSA. Starting from telemetry records and a raster stack of environmental covariates, we will:
\begin{enumerate}
\sphinxsetlistlabels{\arabic}{enumi}{enumii}{}{.}%
\item {} 
\sphinxAtStartPar
prepare and project relocation data;

\item {} 
\sphinxAtStartPar
define the spatial domain considered available;

\item {} 
\sphinxAtStartPar
sample environmental conditions at used and available locations;

\item {} 
\sphinxAtStartPar
inspect predictor distributions;

\item {} 
\sphinxAtStartPar
compare a biologically defined set of candidate models;

\item {} 
\sphinxAtStartPar
fit and project a resource\sphinxhyphen{}selection function;

\item {} 
\sphinxAtStartPar
evaluate transferability among individuals.

\end{enumerate}

\sphinxAtStartPar
The final section briefly shows how the same used\textendash{}available design can be extended to hierarchical Bayesian inference.


\section{Installation}
\label{\detokenize{getting_started:installation}}
\sphinxAtStartPar
Clone the repository and install hrHSA with its geospatial and distributed\sphinxhyphen{}computing dependencies:

\begin{sphinxVerbatim}[commandchars=\\\{\}]
git\PYG{+w}{ }clone\PYG{+w}{ }\PYGZlt{}repository\PYGZhy{}url\PYGZgt{}
\PYG{n+nb}{cd}\PYG{+w}{ }hrHSA
python\PYG{+w}{ }\PYGZhy{}m\PYG{+w}{ }pip\PYG{+w}{ }install\PYG{+w}{ }\PYGZhy{}e\PYG{+w}{ }\PYG{l+s+s2}{\PYGZdq{}.[hpc]\PYGZdq{}}
\end{sphinxVerbatim}

\sphinxAtStartPar
Import the components used throughout this guide:

\begin{sphinxVerbatim}[commandchars=\\\{\}]
\PYG{k+kn}{import}\PYG{+w}{ }\PYG{n+nn}{numpy}\PYG{+w}{ }\PYG{k}{as}\PYG{+w}{ }\PYG{n+nn}{np}
\PYG{k+kn}{import}\PYG{+w}{ }\PYG{n+nn}{pandas}\PYG{+w}{ }\PYG{k}{as}\PYG{+w}{ }\PYG{n+nn}{pd}

\PYG{k+kn}{from}\PYG{+w}{ }\PYG{n+nn}{hsa}\PYG{+w}{ }\PYG{k+kn}{import} \PYG{n}{FeatureSpec}
\PYG{k+kn}{from}\PYG{+w}{ }\PYG{n+nn}{hsa}\PYG{n+nn}{.}\PYG{n+nn}{movement}\PYG{n+nn}{.}\PYG{n+nn}{geometry}\PYG{+w}{ }\PYG{k+kn}{import} \PYG{n}{prepare\PYGZus{}trajectory\PYGZus{}data}
\PYG{k+kn}{from}\PYG{+w}{ }\PYG{n+nn}{hsa}\PYG{n+nn}{.}\PYG{n+nn}{compute}\PYG{+w}{ }\PYG{k+kn}{import} \PYG{n}{open\PYGZus{}raster\PYGZus{}stack\PYGZus{}zarr}\PYG{p}{,} \PYG{n}{suggest\PYGZus{}xy\PYGZus{}chunks}
\PYG{k+kn}{from}\PYG{+w}{ }\PYG{n+nn}{hsa}\PYG{n+nn}{.}\PYG{n+nn}{sampling}\PYG{+w}{ }\PYG{k+kn}{import} \PYG{p}{(}
    \PYG{n}{get\PYGZus{}availability\PYGZus{}domain}\PYG{p}{,}
    \PYG{n}{sample\PYGZus{}available\PYGZus{}points}\PYG{p}{,}
    \PYG{n}{sample\PYGZus{}raster\PYGZus{}stack}\PYG{p}{,}
\PYG{p}{)}
\PYG{k+kn}{from}\PYG{+w}{ }\PYG{n+nn}{hsa}\PYG{n+nn}{.}\PYG{n+nn}{diagnostics}\PYG{+w}{ }\PYG{k+kn}{import} \PYG{n}{inspect\PYGZus{}sampled\PYGZus{}use\PYGZus{}available}
\PYG{k+kn}{from}\PYG{+w}{ }\PYG{n+nn}{hsa}\PYG{n+nn}{.}\PYG{n+nn}{rsf}\PYG{+w}{ }\PYG{k+kn}{import} \PYG{p}{(}
    \PYG{n}{compare\PYGZus{}single\PYGZus{}predictors}\PYG{p}{,}
    \PYG{n}{evaluate\PYGZus{}linear\PYGZus{}candidates\PYGZus{}up\PYGZus{}to\PYGZus{}k}\PYG{p}{,}
    \PYG{n}{fit\PYGZus{}rsf}\PYG{p}{,}
    \PYG{n}{predict\PYGZus{}rsf\PYGZus{}surface}\PYG{p}{,}
\PYG{p}{)}
\PYG{k+kn}{from}\PYG{+w}{ }\PYG{n+nn}{hsa}\PYG{n+nn}{.}\PYG{n+nn}{rsf}\PYG{n+nn}{.}\PYG{n+nn}{cv}\PYG{+w}{ }\PYG{k+kn}{import} \PYG{n}{leave\PYGZus{}one\PYGZus{}individual\PYGZus{}out\PYGZus{}rsf}
\PYG{k+kn}{from}\PYG{+w}{ }\PYG{n+nn}{hsa}\PYG{n+nn}{.}\PYG{n+nn}{rsf}\PYG{n+nn}{.}\PYG{n+nn}{validation}\PYG{+w}{ }\PYG{k+kn}{import} \PYG{n}{plot\PYGZus{}boyce\PYGZus{}curves}
\end{sphinxVerbatim}


\section{1. Prepare relocation data}
\label{\detokenize{getting_started:prepare-relocation-data}}
\sphinxAtStartPar
hrHSA expects timestamps, individual identifiers, and spatial coordinates. Relocations should be represented in a projected coordinate reference system appropriate for the study area.

\begin{sphinxVerbatim}[commandchars=\\\{\}]
\PYG{n}{raw} \PYG{o}{=} \PYG{n}{pd}\PYG{o}{.}\PYG{n}{read\PYGZus{}csv}\PYG{p}{(}\PYG{l+s+s2}{\PYGZdq{}}\PYG{l+s+s2}{data/relocations.csv}\PYG{l+s+s2}{\PYGZdq{}}\PYG{p}{)}

\PYG{n}{reloc} \PYG{o}{=} \PYG{n}{prepare\PYGZus{}trajectory\PYGZus{}data}\PYG{p}{(}
    \PYG{n}{raw}\PYG{p}{,}
    \PYG{n}{id\PYGZus{}col}\PYG{o}{=}\PYG{l+s+s2}{\PYGZdq{}}\PYG{l+s+s2}{Individual\PYGZus{}ID}\PYG{l+s+s2}{\PYGZdq{}}\PYG{p}{,}
    \PYG{n}{timestamp\PYGZus{}col}\PYG{o}{=}\PYG{l+s+s2}{\PYGZdq{}}\PYG{l+s+s2}{Timestamp}\PYG{l+s+s2}{\PYGZdq{}}\PYG{p}{,}
    \PYG{n}{lon\PYGZus{}col}\PYG{o}{=}\PYG{l+s+s2}{\PYGZdq{}}\PYG{l+s+s2}{longitude}\PYG{l+s+s2}{\PYGZdq{}}\PYG{p}{,}
    \PYG{n}{lat\PYGZus{}col}\PYG{o}{=}\PYG{l+s+s2}{\PYGZdq{}}\PYG{l+s+s2}{latitude}\PYG{l+s+s2}{\PYGZdq{}}\PYG{p}{,}
    \PYG{n}{source\PYGZus{}crs}\PYG{o}{=}\PYG{l+s+s2}{\PYGZdq{}}\PYG{l+s+s2}{EPSG:4326}\PYG{l+s+s2}{\PYGZdq{}}\PYG{p}{,}
    \PYG{n}{target\PYGZus{}crs}\PYG{o}{=}\PYG{l+s+s2}{\PYGZdq{}}\PYG{l+s+s2}{EPSG:32733}\PYG{l+s+s2}{\PYGZdq{}}\PYG{p}{,}  \PYG{c+c1}{\PYGZsh{} replace with the CRS of the raster stack}
    \PYG{n}{round\PYGZus{}freq}\PYG{o}{=}\PYG{l+s+s2}{\PYGZdq{}}\PYG{l+s+s2}{h}\PYG{l+s+s2}{\PYGZdq{}}\PYG{p}{,}
    \PYG{n}{drop\PYGZus{}duplicate\PYGZus{}fixes}\PYG{o}{=}\PYG{k+kc}{True}\PYG{p}{,}
\PYG{p}{)}

\PYG{n}{reloc}\PYG{o}{.}\PYG{n}{head}\PYG{p}{(}\PYG{p}{)}
\end{sphinxVerbatim}

\sphinxAtStartPar
The function parses timestamps, removes invalid records, projects coordinates, sorts trajectories, and optionally removes duplicate observations within a specified temporal interval.

\begin{sphinxadmonition}{important}{Important:}
\sphinxAtStartPar
The target CRS must use meaningful planar units and should match the environmental raster stack. Geographic longitude\textendash{}latitude coordinates should not be used for distance\sphinxhyphen{}, area\sphinxhyphen{}, or availability\sphinxhyphen{}based operations.
\end{sphinxadmonition}


\section{2. Open the environmental raster stack}
\label{\detokenize{getting_started:open-the-environmental-raster-stack}}
\sphinxAtStartPar
Environmental covariates can be stored as a Zarr\sphinxhyphen{}backed \sphinxcode{\sphinxupquote{xarray.DataArray}} with dimensions \sphinxcode{\sphinxupquote{band}}, \sphinxcode{\sphinxupquote{y}}, and \sphinxcode{\sphinxupquote{x}}.

\begin{sphinxVerbatim}[commandchars=\\\{\}]
\PYG{n}{env} \PYG{o}{=} \PYG{n}{open\PYGZus{}raster\PYGZus{}stack\PYGZus{}zarr}\PYG{p}{(}
    \PYG{l+s+s2}{\PYGZdq{}}\PYG{l+s+s2}{data/environment.zarr}\PYG{l+s+s2}{\PYGZdq{}}\PYG{p}{,}
    \PYG{n}{name}\PYG{o}{=}\PYG{l+s+s2}{\PYGZdq{}}\PYG{l+s+s2}{env}\PYG{l+s+s2}{\PYGZdq{}}\PYG{p}{,}
    \PYG{n}{chunks}\PYG{o}{=}\PYG{k+kc}{None}\PYG{p}{,}
\PYG{p}{)}

\PYG{n}{chunks} \PYG{o}{=} \PYG{n}{suggest\PYGZus{}xy\PYGZus{}chunks}\PYG{p}{(}
    \PYG{n}{env}\PYG{p}{,}
    \PYG{n}{target\PYGZus{}chunk\PYGZus{}mb}\PYG{o}{=}\PYG{l+m+mi}{128}\PYG{p}{,}
\PYG{p}{)}

\PYG{n}{env} \PYG{o}{=} \PYG{n}{env}\PYG{o}{.}\PYG{n}{chunk}\PYG{p}{(}\PYG{n}{chunks}\PYG{p}{)}

\PYG{n+nb}{print}\PYG{p}{(}\PYG{n}{chunks}\PYG{p}{)}
\PYG{n+nb}{print}\PYG{p}{(}\PYG{n}{env}\PYG{p}{)}
\end{sphinxVerbatim}

\sphinxAtStartPar
\sphinxcode{\sphinxupquote{suggest\_xy\_chunks}} chooses approximately square spatial chunks while retaining all predictor bands within each chunk. This provides a defensible starting point for raster sampling and prediction, although chunk sizes should ultimately be benchmarked against the available memory, storage system, and workload.

\sphinxAtStartPar
Confirm that the relocation and raster coordinate systems agree:

\begin{sphinxVerbatim}[commandchars=\\\{\}]
\PYG{k}{assert} \PYG{n}{reloc}\PYG{o}{.}\PYG{n}{crs} \PYG{o}{==} \PYG{n}{env}\PYG{o}{.}\PYG{n}{rio}\PYG{o}{.}\PYG{n}{crs}
\end{sphinxVerbatim}


\section{3. Define availability}
\label{\detokenize{getting_started:define-availability}}
\sphinxAtStartPar
A resource\sphinxhyphen{}selection analysis compares observed use with environmental conditions that were accessible but not necessarily used. The definition of availability is therefore part of the ecological model rather than merely a computational preprocessing decision.

\sphinxAtStartPar
Here, separate availability domains are estimated for each individual using a centroid\sphinxhyphen{}trimmed 95\% minimum convex polygon:

\begin{sphinxVerbatim}[commandchars=\\\{\}]
\PYG{n}{domains} \PYG{o}{=} \PYG{n}{get\PYGZus{}availability\PYGZus{}domain}\PYG{p}{(}
    \PYG{n}{reloc}\PYG{p}{,}
    \PYG{n}{id\PYGZus{}col}\PYG{o}{=}\PYG{l+s+s2}{\PYGZdq{}}\PYG{l+s+s2}{Individual\PYGZus{}ID}\PYG{l+s+s2}{\PYGZdq{}}\PYG{p}{,}
    \PYG{n}{estimator}\PYG{o}{=}\PYG{l+s+s2}{\PYGZdq{}}\PYG{l+s+s2}{mcp}\PYG{l+s+s2}{\PYGZdq{}}\PYG{p}{,}
    \PYG{n}{quantile}\PYG{o}{=}\PYG{l+m+mf}{0.95}\PYG{p}{,}
    \PYG{n}{min\PYGZus{}points}\PYG{o}{=}\PYG{l+m+mi}{5}\PYG{p}{,}
\PYG{p}{)}

\PYG{n}{domains}\PYG{o}{.}\PYG{n}{head}\PYG{p}{(}\PYG{p}{)}
\end{sphinxVerbatim}

\sphinxAtStartPar
Available locations are then sampled independently within each individual domain:

\begin{sphinxVerbatim}[commandchars=\\\{\}]
\PYG{n}{samples} \PYG{o}{=} \PYG{n}{sample\PYGZus{}available\PYGZus{}points}\PYG{p}{(}
    \PYG{n}{domains}\PYG{p}{,}
    \PYG{n}{used}\PYG{o}{=}\PYG{n}{reloc}\PYG{p}{,}
    \PYG{n}{id\PYGZus{}col}\PYG{o}{=}\PYG{l+s+s2}{\PYGZdq{}}\PYG{l+s+s2}{Individual\PYGZus{}ID}\PYG{l+s+s2}{\PYGZdq{}}\PYG{p}{,}
    \PYG{n}{n\PYGZus{}per\PYGZus{}used}\PYG{o}{=}\PYG{l+m+mi}{10}\PYG{p}{,}
    \PYG{n}{seed}\PYG{o}{=}\PYG{l+m+mi}{42}\PYG{p}{,}
\PYG{p}{)}

\PYG{n}{samples}\PYG{p}{[}\PYG{l+s+s2}{\PYGZdq{}}\PYG{l+s+s2}{used}\PYG{l+s+s2}{\PYGZdq{}}\PYG{p}{]}\PYG{o}{.}\PYG{n}{value\PYGZus{}counts}\PYG{p}{(}\PYG{p}{)}
\end{sphinxVerbatim}

\sphinxAtStartPar
The resulting table contains both observed relocations (\sphinxcode{\sphinxupquote{used=True}}) and sampled available locations (\sphinxcode{\sphinxupquote{used=False}}).

\begin{sphinxadmonition}{note}{Note:}
\sphinxAtStartPar
The number of available points controls the numerical approximation of the available environment. It does not represent the number of locations an animal could literally have visited.
\end{sphinxadmonition}


\section{4. Extract environmental covariates}
\label{\detokenize{getting_started:extract-environmental-covariates}}
\sphinxAtStartPar
Sample all raster bands at the used and available locations:

\begin{sphinxVerbatim}[commandchars=\\\{\}]
\PYG{n}{sampled} \PYG{o}{=} \PYG{n}{sample\PYGZus{}raster\PYGZus{}stack}\PYG{p}{(}
    \PYG{n}{samples}\PYG{p}{,}
    \PYG{n}{env}\PYG{p}{,}
    \PYG{n}{id\PYGZus{}cols}\PYG{o}{=}\PYG{l+s+s2}{\PYGZdq{}}\PYG{l+s+s2}{Individual\PYGZus{}ID}\PYG{l+s+s2}{\PYGZdq{}}\PYG{p}{,}
\PYG{p}{)}

\PYG{n}{sampled}\PYG{o}{.}\PYG{n}{head}\PYG{p}{(}\PYG{p}{)}
\end{sphinxVerbatim}

\sphinxAtStartPar
Each row now represents one used or available location together with its environmental covariates.


\section{5. Inspect used and available environments}
\label{\detokenize{getting_started:inspect-used-and-available-environments}}
\sphinxAtStartPar
Before fitting a model, inspect missingness, predictor distributions, and the degree of separation between used and available samples.

\begin{sphinxVerbatim}[commandchars=\\\{\}]
\PYG{n}{inspection} \PYG{o}{=} \PYG{n}{inspect\PYGZus{}sampled\PYGZus{}use\PYGZus{}available}\PYG{p}{(}
    \PYG{n}{sampled}\PYG{p}{,}
    \PYG{n}{id\PYGZus{}col}\PYG{o}{=}\PYG{l+s+s2}{\PYGZdq{}}\PYG{l+s+s2}{Individual\PYGZus{}ID}\PYG{l+s+s2}{\PYGZdq{}}\PYG{p}{,}
    \PYG{n}{plot}\PYG{o}{=}\PYG{k+kc}{True}\PYG{p}{,}
    \PYG{n}{plot\PYGZus{}top\PYGZus{}n}\PYG{o}{=}\PYG{l+m+mi}{8}\PYG{p}{,}
\PYG{p}{)}

\PYG{n}{inspection}\PYG{p}{[}\PYG{l+s+s2}{\PYGZdq{}}\PYG{l+s+s2}{pooled\PYGZus{}stats}\PYG{l+s+s2}{\PYGZdq{}}\PYG{p}{]}
\end{sphinxVerbatim}

\sphinxAtStartPar
The diagnostic summary includes distributional differences such as the Kolmogorov\textendash{}Smirnov statistic, Wasserstein distance, differences in means and medians, and the common\sphinxhyphen{}language effect size.

\sphinxAtStartPar
These summaries are descriptive. They are useful for identifying data problems, implausible contrasts, and potentially informative predictors, but they do not replace a defined ecological hypothesis or multivariable model.


\section{6. Evaluate candidate models}
\label{\detokenize{getting_started:evaluate-candidate-models}}
\sphinxAtStartPar
Begin with a biologically justified set of candidate predictors:

\begin{sphinxVerbatim}[commandchars=\\\{\}]
\PYG{n}{candidate\PYGZus{}predictors} \PYG{o}{=} \PYG{p}{[}
    \PYG{l+s+s2}{\PYGZdq{}}\PYG{l+s+s2}{elevation}\PYG{l+s+s2}{\PYGZdq{}}\PYG{p}{,}
    \PYG{l+s+s2}{\PYGZdq{}}\PYG{l+s+s2}{slope}\PYG{l+s+s2}{\PYGZdq{}}\PYG{p}{,}
    \PYG{l+s+s2}{\PYGZdq{}}\PYG{l+s+s2}{ndvi}\PYG{l+s+s2}{\PYGZdq{}}\PYG{p}{,}
    \PYG{l+s+s2}{\PYGZdq{}}\PYG{l+s+s2}{distance\PYGZus{}to\PYGZus{}water}\PYG{l+s+s2}{\PYGZdq{}}\PYG{p}{,}
\PYG{p}{]}
\end{sphinxVerbatim}

\sphinxAtStartPar
Compare single\sphinxhyphen{}predictor models:

\begin{sphinxVerbatim}[commandchars=\\\{\}]
\PYG{n}{single\PYGZus{}predictor\PYGZus{}models} \PYG{o}{=} \PYG{n}{compare\PYGZus{}single\PYGZus{}predictors}\PYG{p}{(}
    \PYG{n}{sampled}\PYG{p}{,}
    \PYG{n}{predictors}\PYG{o}{=}\PYG{n}{candidate\PYGZus{}predictors}\PYG{p}{,}
\PYG{p}{)}

\PYG{n}{single\PYGZus{}predictor\PYGZus{}models}\PYG{p}{[}
    \PYG{p}{[}
        \PYG{l+s+s2}{\PYGZdq{}}\PYG{l+s+s2}{predictor}\PYG{l+s+s2}{\PYGZdq{}}\PYG{p}{,}
        \PYG{l+s+s2}{\PYGZdq{}}\PYG{l+s+s2}{AIC}\PYG{l+s+s2}{\PYGZdq{}}\PYG{p}{,}
        \PYG{l+s+s2}{\PYGZdq{}}\PYG{l+s+s2}{delta\PYGZus{}aic}\PYG{l+s+s2}{\PYGZdq{}}\PYG{p}{,}
        \PYG{l+s+s2}{\PYGZdq{}}\PYG{l+s+s2}{akaike\PYGZus{}weight}\PYG{l+s+s2}{\PYGZdq{}}\PYG{p}{,}
        \PYG{l+s+s2}{\PYGZdq{}}\PYG{l+s+s2}{pseudo\PYGZus{}r2}\PYG{l+s+s2}{\PYGZdq{}}\PYG{p}{,}
        \PYG{l+s+s2}{\PYGZdq{}}\PYG{l+s+s2}{converged}\PYG{l+s+s2}{\PYGZdq{}}\PYG{p}{,}
    \PYG{p}{]}
\PYG{p}{]}
\end{sphinxVerbatim}

\sphinxAtStartPar
Small candidate combinations can then be evaluated:

\begin{sphinxVerbatim}[commandchars=\\\{\}]
\PYG{n}{candidate\PYGZus{}models} \PYG{o}{=} \PYG{n}{evaluate\PYGZus{}linear\PYGZus{}candidates\PYGZus{}up\PYGZus{}to\PYGZus{}k}\PYG{p}{(}
    \PYG{n}{sampled}\PYG{p}{,}
    \PYG{n}{predictor\PYGZus{}cols}\PYG{o}{=}\PYG{n}{candidate\PYGZus{}predictors}\PYG{p}{,}
    \PYG{n}{max\PYGZus{}k}\PYG{o}{=}\PYG{l+m+mi}{3}\PYG{p}{,}
    \PYG{n}{n\PYGZus{}jobs}\PYG{o}{=}\PYG{l+m+mi}{4}\PYG{p}{,}
\PYG{p}{)}

\PYG{n}{candidate\PYGZus{}models}\PYG{o}{.}\PYG{n}{head}\PYG{p}{(}\PYG{l+m+mi}{10}\PYG{p}{)}
\end{sphinxVerbatim}

\begin{sphinxadmonition}{warning}{Warning:}
\sphinxAtStartPar
The number of possible models grows combinatorially. Candidate predictors, transformations, and interactions should be constrained by ecological reasoning before model fitting. AIC ranking quantifies relative support within the supplied candidate set; it does not discover causal structure automatically.
\end{sphinxadmonition}


\section{7. Specify and fit the RSF}
\label{\detokenize{getting_started:specify-and-fit-the-rsf}}
\sphinxAtStartPar
\sphinxcode{\sphinxupquote{FeatureSpec}} defines the model matrix explicitly. Continuous predictors are standardized during fitting, while quadratic terms and interactions are constructed on the standardized scale.

\begin{sphinxVerbatim}[commandchars=\\\{\}]
\PYG{n}{spec} \PYG{o}{=} \PYG{n}{FeatureSpec}\PYG{p}{(}
    \PYG{n}{linear}\PYG{o}{=}\PYG{p}{[}
        \PYG{l+s+s2}{\PYGZdq{}}\PYG{l+s+s2}{distance\PYGZus{}to\PYGZus{}water}\PYG{l+s+s2}{\PYGZdq{}}\PYG{p}{,}
        \PYG{l+s+s2}{\PYGZdq{}}\PYG{l+s+s2}{ndvi}\PYG{l+s+s2}{\PYGZdq{}}\PYG{p}{,}
        \PYG{l+s+s2}{\PYGZdq{}}\PYG{l+s+s2}{slope}\PYG{l+s+s2}{\PYGZdq{}}\PYG{p}{,}
    \PYG{p}{]}\PYG{p}{,}
    \PYG{n}{quadratic}\PYG{o}{=}\PYG{p}{[}
        \PYG{l+s+s2}{\PYGZdq{}}\PYG{l+s+s2}{distance\PYGZus{}to\PYGZus{}water}\PYG{l+s+s2}{\PYGZdq{}}\PYG{p}{,}
    \PYG{p}{]}\PYG{p}{,}
    \PYG{n}{interactions}\PYG{o}{=}\PYG{p}{[}
        \PYG{p}{(}\PYG{l+s+s2}{\PYGZdq{}}\PYG{l+s+s2}{ndvi}\PYG{l+s+s2}{\PYGZdq{}}\PYG{p}{,} \PYG{l+s+s2}{\PYGZdq{}}\PYG{l+s+s2}{slope}\PYG{l+s+s2}{\PYGZdq{}}\PYG{p}{)}\PYG{p}{,}
    \PYG{p}{]}\PYG{p}{,}
    \PYG{n}{add\PYGZus{}const}\PYG{o}{=}\PYG{k+kc}{True}\PYG{p}{,}
\PYG{p}{)}
\end{sphinxVerbatim}

\sphinxAtStartPar
Fit the used\textendash{}available logistic regression:

\begin{sphinxVerbatim}[commandchars=\\\{\}]
\PYG{n}{model}\PYG{p}{,} \PYG{n}{scaler}\PYG{p}{,} \PYG{n}{fitted\PYGZus{}spec}\PYG{p}{,} \PYG{n}{meta} \PYG{o}{=} \PYG{n}{fit\PYGZus{}rsf}\PYG{p}{(}
    \PYG{n}{sampled}\PYG{p}{,}
    \PYG{n}{spec}\PYG{p}{,}
\PYG{p}{)}

\PYG{n+nb}{print}\PYG{p}{(}\PYG{n}{model}\PYG{o}{.}\PYG{n}{summary}\PYG{p}{(}\PYG{p}{)}\PYG{p}{)}
\end{sphinxVerbatim}

\sphinxAtStartPar
For continuous predictors, positive coefficients indicate increasing relative selection with increasing predictor values, conditional on the remaining terms. Negative coefficients indicate decreasing relative selection.

\sphinxAtStartPar
Because the model is fitted to a researcher\sphinxhyphen{}defined ratio of used and available locations, its intercept and raw logistic probabilities generally do not represent absolute occurrence probabilities. The primary inferential quantities are relative selection coefficients and contrasts in relative selection strength.


\section{8. Predict the selection surface}
\label{\detokenize{getting_started:predict-the-selection-surface}}
\sphinxAtStartPar
Project the fitted model over the environmental raster stack:

\begin{sphinxVerbatim}[commandchars=\\\{\}]
\PYG{n}{rsf} \PYG{o}{=} \PYG{n}{predict\PYGZus{}rsf\PYGZus{}surface}\PYG{p}{(}
    \PYG{n}{env}\PYG{p}{,}
    \PYG{n}{model}\PYG{p}{,}
    \PYG{n}{scaler}\PYG{p}{,}
    \PYG{n}{fitted\PYGZus{}spec}\PYG{p}{,}
    \PYG{n}{meta}\PYG{p}{,}
\PYG{p}{)}

\PYG{n}{rsf}
\end{sphinxVerbatim}

\sphinxAtStartPar
The output is an \sphinxcode{\sphinxupquote{xarray.DataArray}} containing the relative selection surface:

\begin{sphinxVerbatim}[commandchars=\\\{\}]
\PYG{n}{rsf}\PYG{o}{.}\PYG{n}{sel}\PYG{p}{(}\PYG{n}{band}\PYG{o}{=}\PYG{l+s+s2}{\PYGZdq{}}\PYG{l+s+s2}{rsf}\PYG{l+s+s2}{\PYGZdq{}}\PYG{p}{)}\PYG{o}{.}\PYG{n}{plot}\PYG{p}{(}
    \PYG{n}{robust}\PYG{o}{=}\PYG{k+kc}{True}\PYG{p}{,}
    \PYG{n}{figsize}\PYG{o}{=}\PYG{p}{(}\PYG{l+m+mi}{9}\PYG{p}{,} \PYG{l+m+mi}{7}\PYG{p}{)}\PYG{p}{,}
\PYG{p}{)}
\end{sphinxVerbatim}

\sphinxAtStartPar
The surface is proportional to

\sphinxAtStartPar
{[}
w(\textbackslash{}mathbf\{x\}) = \textbackslash{}exp!\textbackslash{}left(
\textbackslash{}beta\_0 + \textbackslash{}boldsymbol\{\textbackslash{}beta\}\textasciicircum{}\{\textbackslash{}mathsf T\}\textbackslash{}mathbf\{x\}
\textbackslash{}right),
{]}

\sphinxAtStartPar
and should therefore be interpreted comparatively: a location with a score of two is estimated to have twice the relative selection strength of a location with a score of one, under the fitted model and availability design.


\section{9. Validate among individuals}
\label{\detokenize{getting_started:validate-among-individuals}}
\sphinxAtStartPar
Randomly splitting autocorrelated relocations can produce optimistic validation results. For multi\sphinxhyphen{}individual datasets, hrHSA therefore supports leave\sphinxhyphen{}one\sphinxhyphen{}individual\sphinxhyphen{}out validation.

\begin{sphinxVerbatim}[commandchars=\\\{\}]
\PYG{p}{(}
    \PYG{n}{cv\PYGZus{}summary}\PYG{p}{,}
    \PYG{n}{cv\PYGZus{}parameters}\PYG{p}{,}
    \PYG{n}{boyce\PYGZus{}bins}\PYG{p}{,}
    \PYG{n}{calibration\PYGZus{}bins}\PYG{p}{,}
    \PYG{n}{cv\PYGZus{}diagnostics}\PYG{p}{,}
\PYG{p}{)} \PYG{o}{=} \PYG{n}{leave\PYGZus{}one\PYGZus{}individual\PYGZus{}out\PYGZus{}rsf}\PYG{p}{(}
    \PYG{n}{reloc}\PYG{p}{,}
    \PYG{n}{env}\PYG{p}{,}
    \PYG{n}{fitted\PYGZus{}spec}\PYG{p}{,}
    \PYG{n}{id\PYGZus{}col}\PYG{o}{=}\PYG{l+s+s2}{\PYGZdq{}}\PYG{l+s+s2}{Individual\PYGZus{}ID}\PYG{l+s+s2}{\PYGZdq{}}\PYG{p}{,}
    \PYG{n}{heldout}\PYG{o}{=}\PYG{l+s+s2}{\PYGZdq{}}\PYG{l+s+s2}{all}\PYG{l+s+s2}{\PYGZdq{}}\PYG{p}{,}
    \PYG{n}{domain\PYGZus{}quantile}\PYG{o}{=}\PYG{l+m+mf}{0.95}\PYG{p}{,}
    \PYG{n}{thin\PYGZus{}train\PYGZus{}dt}\PYG{o}{=}\PYG{l+s+s2}{\PYGZdq{}}\PYG{l+s+s2}{12h}\PYG{l+s+s2}{\PYGZdq{}}\PYG{p}{,}
    \PYG{n}{sampling\PYGZus{}factor\PYGZus{}train}\PYG{o}{=}\PYG{l+m+mi}{10}\PYG{p}{,}
    \PYG{n}{n\PYGZus{}background\PYGZus{}boyce}\PYG{o}{=}\PYG{l+m+mi}{100\PYGZus{}000}\PYG{p}{,}
    \PYG{n}{seed}\PYG{o}{=}\PYG{l+m+mi}{42}\PYG{p}{,}
\PYG{p}{)}

\PYG{n}{cv\PYGZus{}summary}
\end{sphinxVerbatim}

\sphinxAtStartPar
In each fold, the model is trained on all but one individual and evaluated against the relocations and availability domain of the held\sphinxhyphen{}out individual. This tests whether population\sphinxhyphen{}level selection relationships transfer to animals not used during model fitting.

\sphinxAtStartPar
Plot the fold\sphinxhyphen{}specific Boyce curves:

\begin{sphinxVerbatim}[commandchars=\\\{\}]
\PYG{n}{fig}\PYG{p}{,} \PYG{n}{ax} \PYG{o}{=} \PYG{n}{plot\PYGZus{}boyce\PYGZus{}curves}\PYG{p}{(}
    \PYG{n}{boyce\PYGZus{}bins}\PYG{p}{,}
    \PYG{n}{id\PYGZus{}col}\PYG{o}{=}\PYG{l+s+s2}{\PYGZdq{}}\PYG{l+s+s2}{heldout\PYGZus{}ID}\PYG{l+s+s2}{\PYGZdq{}}\PYG{p}{,}
\PYG{p}{)}
\end{sphinxVerbatim}

\sphinxAtStartPar
The Boyce diagnostic evaluates whether observed locations become progressively more frequent than expected under random use as predicted relative selection increases. It should be interpreted alongside fold\sphinxhyphen{}specific coefficients, calibration diagnostics, sample sizes, and the ecological structure of the held\sphinxhyphen{}out individuals.


\section{A hierarchical Bayesian extension}
\label{\detokenize{getting_started:a-hierarchical-bayesian-extension}}
\sphinxAtStartPar
The frequentist workflow estimates a common population\sphinxhyphen{}level selection relationship. When repeated observations from multiple individuals are available, a hierarchical model can estimate both population\sphinxhyphen{}level effects and among\sphinxhyphen{}individual variation.

\sphinxAtStartPar
The following abbreviated PyMC model gives each individual its own intercept and slope while partially pooling those effects toward the population mean.

\begin{sphinxVerbatim}[commandchars=\\\{\}]
\PYG{k+kn}{import}\PYG{+w}{ }\PYG{n+nn}{pymc}\PYG{+w}{ }\PYG{k}{as}\PYG{+w}{ }\PYG{n+nn}{pm}

\PYG{n}{predictor} \PYG{o}{=} \PYG{l+s+s2}{\PYGZdq{}}\PYG{l+s+s2}{distance\PYGZus{}to\PYGZus{}water}\PYG{l+s+s2}{\PYGZdq{}}

\PYG{n}{bayes\PYGZus{}df} \PYG{o}{=} \PYG{n}{sampled}\PYG{p}{[}
    \PYG{p}{[}\PYG{l+s+s2}{\PYGZdq{}}\PYG{l+s+s2}{used}\PYG{l+s+s2}{\PYGZdq{}}\PYG{p}{,} \PYG{l+s+s2}{\PYGZdq{}}\PYG{l+s+s2}{Individual\PYGZus{}ID}\PYG{l+s+s2}{\PYGZdq{}}\PYG{p}{,} \PYG{n}{predictor}\PYG{p}{]}
\PYG{p}{]}\PYG{o}{.}\PYG{n}{dropna}\PYG{p}{(}\PYG{p}{)}\PYG{o}{.}\PYG{n}{copy}\PYG{p}{(}\PYG{p}{)}

\PYG{n}{x} \PYG{o}{=} \PYG{n}{bayes\PYGZus{}df}\PYG{p}{[}\PYG{n}{predictor}\PYG{p}{]}\PYG{o}{.}\PYG{n}{to\PYGZus{}numpy}\PYG{p}{(}\PYG{n}{dtype}\PYG{o}{=}\PYG{n+nb}{float}\PYG{p}{)}
\PYG{n}{x} \PYG{o}{=} \PYG{p}{(}\PYG{n}{x} \PYG{o}{\PYGZhy{}} \PYG{n}{x}\PYG{o}{.}\PYG{n}{mean}\PYG{p}{(}\PYG{p}{)}\PYG{p}{)} \PYG{o}{/} \PYG{n}{x}\PYG{o}{.}\PYG{n}{std}\PYG{p}{(}\PYG{p}{)}

\PYG{n}{y} \PYG{o}{=} \PYG{n}{bayes\PYGZus{}df}\PYG{p}{[}\PYG{l+s+s2}{\PYGZdq{}}\PYG{l+s+s2}{used}\PYG{l+s+s2}{\PYGZdq{}}\PYG{p}{]}\PYG{o}{.}\PYG{n}{to\PYGZus{}numpy}\PYG{p}{(}\PYG{n}{dtype}\PYG{o}{=}\PYG{n+nb}{int}\PYG{p}{)}

\PYG{n}{id\PYGZus{}idx}\PYG{p}{,} \PYG{n}{individual\PYGZus{}names} \PYG{o}{=} \PYG{n}{pd}\PYG{o}{.}\PYG{n}{factorize}\PYG{p}{(}
    \PYG{n}{bayes\PYGZus{}df}\PYG{p}{[}\PYG{l+s+s2}{\PYGZdq{}}\PYG{l+s+s2}{Individual\PYGZus{}ID}\PYG{l+s+s2}{\PYGZdq{}}\PYG{p}{]}\PYG{p}{,}
    \PYG{n}{sort}\PYG{o}{=}\PYG{k+kc}{True}\PYG{p}{,}
\PYG{p}{)}

\PYG{n}{coords} \PYG{o}{=} \PYG{p}{\PYGZob{}}
    \PYG{l+s+s2}{\PYGZdq{}}\PYG{l+s+s2}{obs}\PYG{l+s+s2}{\PYGZdq{}}\PYG{p}{:} \PYG{n}{np}\PYG{o}{.}\PYG{n}{arange}\PYG{p}{(}\PYG{n+nb}{len}\PYG{p}{(}\PYG{n}{bayes\PYGZus{}df}\PYG{p}{)}\PYG{p}{)}\PYG{p}{,}
    \PYG{l+s+s2}{\PYGZdq{}}\PYG{l+s+s2}{individual}\PYG{l+s+s2}{\PYGZdq{}}\PYG{p}{:} \PYG{n}{individual\PYGZus{}names}\PYG{o}{.}\PYG{n}{tolist}\PYG{p}{(}\PYG{p}{)}\PYG{p}{,}
\PYG{p}{\PYGZcb{}}
\end{sphinxVerbatim}

\begin{sphinxVerbatim}[commandchars=\\\{\}]
\PYG{k}{with} \PYG{n}{pm}\PYG{o}{.}\PYG{n}{Model}\PYG{p}{(}\PYG{n}{coords}\PYG{o}{=}\PYG{n}{coords}\PYG{p}{)} \PYG{k}{as} \PYG{n}{hierarchical\PYGZus{}rsf}\PYG{p}{:}

    \PYG{c+c1}{\PYGZsh{} Data containers}
    \PYG{n}{x\PYGZus{}data} \PYG{o}{=} \PYG{n}{pm}\PYG{o}{.}\PYG{n}{Data}\PYG{p}{(}\PYG{l+s+s2}{\PYGZdq{}}\PYG{l+s+s2}{x}\PYG{l+s+s2}{\PYGZdq{}}\PYG{p}{,} \PYG{n}{x}\PYG{p}{,} \PYG{n}{dims}\PYG{o}{=}\PYG{l+s+s2}{\PYGZdq{}}\PYG{l+s+s2}{obs}\PYG{l+s+s2}{\PYGZdq{}}\PYG{p}{)}
    \PYG{n}{id\PYGZus{}data} \PYG{o}{=} \PYG{n}{pm}\PYG{o}{.}\PYG{n}{Data}\PYG{p}{(}\PYG{l+s+s2}{\PYGZdq{}}\PYG{l+s+s2}{id\PYGZus{}idx}\PYG{l+s+s2}{\PYGZdq{}}\PYG{p}{,} \PYG{n}{id\PYGZus{}idx}\PYG{p}{,} \PYG{n}{dims}\PYG{o}{=}\PYG{l+s+s2}{\PYGZdq{}}\PYG{l+s+s2}{obs}\PYG{l+s+s2}{\PYGZdq{}}\PYG{p}{)}
    \PYG{n}{y\PYGZus{}data} \PYG{o}{=} \PYG{n}{pm}\PYG{o}{.}\PYG{n}{Data}\PYG{p}{(}\PYG{l+s+s2}{\PYGZdq{}}\PYG{l+s+s2}{y}\PYG{l+s+s2}{\PYGZdq{}}\PYG{p}{,} \PYG{n}{y}\PYG{p}{,} \PYG{n}{dims}\PYG{o}{=}\PYG{l+s+s2}{\PYGZdq{}}\PYG{l+s+s2}{obs}\PYG{l+s+s2}{\PYGZdq{}}\PYG{p}{)}

    \PYG{c+c1}{\PYGZsh{} Population\PYGZhy{}level intercept and slope}
    \PYG{n}{alpha} \PYG{o}{=} \PYG{n}{pm}\PYG{o}{.}\PYG{n}{Normal}\PYG{p}{(}\PYG{l+s+s2}{\PYGZdq{}}\PYG{l+s+s2}{alpha}\PYG{l+s+s2}{\PYGZdq{}}\PYG{p}{,} \PYG{n}{mu}\PYG{o}{=}\PYG{l+m+mf}{0.0}\PYG{p}{,} \PYG{n}{sigma}\PYG{o}{=}\PYG{l+m+mf}{2.5}\PYG{p}{)}
    \PYG{n}{beta} \PYG{o}{=} \PYG{n}{pm}\PYG{o}{.}\PYG{n}{Normal}\PYG{p}{(}\PYG{l+s+s2}{\PYGZdq{}}\PYG{l+s+s2}{beta}\PYG{l+s+s2}{\PYGZdq{}}\PYG{p}{,} \PYG{n}{mu}\PYG{o}{=}\PYG{l+m+mf}{0.0}\PYG{p}{,} \PYG{n}{sigma}\PYG{o}{=}\PYG{l+m+mf}{1.0}\PYG{p}{)}

    \PYG{c+c1}{\PYGZsh{} Among\PYGZhy{}individual variation}
    \PYG{n}{sigma\PYGZus{}alpha} \PYG{o}{=} \PYG{n}{pm}\PYG{o}{.}\PYG{n}{Exponential}\PYG{p}{(}\PYG{l+s+s2}{\PYGZdq{}}\PYG{l+s+s2}{sigma\PYGZus{}alpha}\PYG{l+s+s2}{\PYGZdq{}}\PYG{p}{,} \PYG{n}{lam}\PYG{o}{=}\PYG{l+m+mf}{1.0}\PYG{p}{)}
    \PYG{n}{sigma\PYGZus{}beta} \PYG{o}{=} \PYG{n}{pm}\PYG{o}{.}\PYG{n}{Exponential}\PYG{p}{(}\PYG{l+s+s2}{\PYGZdq{}}\PYG{l+s+s2}{sigma\PYGZus{}beta}\PYG{l+s+s2}{\PYGZdq{}}\PYG{p}{,} \PYG{n}{lam}\PYG{o}{=}\PYG{l+m+mf}{1.0}\PYG{p}{)}

    \PYG{c+c1}{\PYGZsh{} Standardized individual deviations}
    \PYG{n}{z\PYGZus{}alpha} \PYG{o}{=} \PYG{n}{pm}\PYG{o}{.}\PYG{n}{Normal}\PYG{p}{(}
        \PYG{l+s+s2}{\PYGZdq{}}\PYG{l+s+s2}{z\PYGZus{}alpha}\PYG{l+s+s2}{\PYGZdq{}}\PYG{p}{,}
        \PYG{n}{mu}\PYG{o}{=}\PYG{l+m+mf}{0.0}\PYG{p}{,}
        \PYG{n}{sigma}\PYG{o}{=}\PYG{l+m+mf}{1.0}\PYG{p}{,}
        \PYG{n}{dims}\PYG{o}{=}\PYG{l+s+s2}{\PYGZdq{}}\PYG{l+s+s2}{individual}\PYG{l+s+s2}{\PYGZdq{}}\PYG{p}{,}
    \PYG{p}{)}
    \PYG{n}{z\PYGZus{}beta} \PYG{o}{=} \PYG{n}{pm}\PYG{o}{.}\PYG{n}{Normal}\PYG{p}{(}
        \PYG{l+s+s2}{\PYGZdq{}}\PYG{l+s+s2}{z\PYGZus{}beta}\PYG{l+s+s2}{\PYGZdq{}}\PYG{p}{,}
        \PYG{n}{mu}\PYG{o}{=}\PYG{l+m+mf}{0.0}\PYG{p}{,}
        \PYG{n}{sigma}\PYG{o}{=}\PYG{l+m+mf}{1.0}\PYG{p}{,}
        \PYG{n}{dims}\PYG{o}{=}\PYG{l+s+s2}{\PYGZdq{}}\PYG{l+s+s2}{individual}\PYG{l+s+s2}{\PYGZdq{}}\PYG{p}{,}
    \PYG{p}{)}

    \PYG{c+c1}{\PYGZsh{} Non\PYGZhy{}centered random effects}
    \PYG{n}{alpha\PYGZus{}ind} \PYG{o}{=} \PYG{n}{pm}\PYG{o}{.}\PYG{n}{Deterministic}\PYG{p}{(}
        \PYG{l+s+s2}{\PYGZdq{}}\PYG{l+s+s2}{alpha\PYGZus{}ind}\PYG{l+s+s2}{\PYGZdq{}}\PYG{p}{,}
        \PYG{n}{z\PYGZus{}alpha} \PYG{o}{*} \PYG{n}{sigma\PYGZus{}alpha}\PYG{p}{,}
        \PYG{n}{dims}\PYG{o}{=}\PYG{l+s+s2}{\PYGZdq{}}\PYG{l+s+s2}{individual}\PYG{l+s+s2}{\PYGZdq{}}\PYG{p}{,}
    \PYG{p}{)}
    \PYG{n}{beta\PYGZus{}ind} \PYG{o}{=} \PYG{n}{pm}\PYG{o}{.}\PYG{n}{Deterministic}\PYG{p}{(}
        \PYG{l+s+s2}{\PYGZdq{}}\PYG{l+s+s2}{beta\PYGZus{}ind}\PYG{l+s+s2}{\PYGZdq{}}\PYG{p}{,}
        \PYG{n}{z\PYGZus{}beta} \PYG{o}{*} \PYG{n}{sigma\PYGZus{}beta}\PYG{p}{,}
        \PYG{n}{dims}\PYG{o}{=}\PYG{l+s+s2}{\PYGZdq{}}\PYG{l+s+s2}{individual}\PYG{l+s+s2}{\PYGZdq{}}\PYG{p}{,}
    \PYG{p}{)}

    \PYG{c+c1}{\PYGZsh{} Total individual effects}
    \PYG{n}{alpha\PYGZus{}total} \PYG{o}{=} \PYG{n}{pm}\PYG{o}{.}\PYG{n}{Deterministic}\PYG{p}{(}
        \PYG{l+s+s2}{\PYGZdq{}}\PYG{l+s+s2}{alpha\PYGZus{}total}\PYG{l+s+s2}{\PYGZdq{}}\PYG{p}{,}
        \PYG{n}{alpha} \PYG{o}{+} \PYG{n}{alpha\PYGZus{}ind}\PYG{p}{,}
        \PYG{n}{dims}\PYG{o}{=}\PYG{l+s+s2}{\PYGZdq{}}\PYG{l+s+s2}{individual}\PYG{l+s+s2}{\PYGZdq{}}\PYG{p}{,}
    \PYG{p}{)}
    \PYG{n}{beta\PYGZus{}total} \PYG{o}{=} \PYG{n}{pm}\PYG{o}{.}\PYG{n}{Deterministic}\PYG{p}{(}
        \PYG{l+s+s2}{\PYGZdq{}}\PYG{l+s+s2}{beta\PYGZus{}total}\PYG{l+s+s2}{\PYGZdq{}}\PYG{p}{,}
        \PYG{n}{beta} \PYG{o}{+} \PYG{n}{beta\PYGZus{}ind}\PYG{p}{,}
        \PYG{n}{dims}\PYG{o}{=}\PYG{l+s+s2}{\PYGZdq{}}\PYG{l+s+s2}{individual}\PYG{l+s+s2}{\PYGZdq{}}\PYG{p}{,}
    \PYG{p}{)}

    \PYG{c+c1}{\PYGZsh{} Linear predictor}
    \PYG{n}{eta} \PYG{o}{=} \PYG{p}{(}
        \PYG{n}{alpha\PYGZus{}total}\PYG{p}{[}\PYG{n}{id\PYGZus{}data}\PYG{p}{]}
        \PYG{o}{+} \PYG{n}{beta\PYGZus{}total}\PYG{p}{[}\PYG{n}{id\PYGZus{}data}\PYG{p}{]} \PYG{o}{*} \PYG{n}{x\PYGZus{}data}
    \PYG{p}{)}

    \PYG{c+c1}{\PYGZsh{} Used\textendash{}available likelihood}
    \PYG{n}{used} \PYG{o}{=} \PYG{n}{pm}\PYG{o}{.}\PYG{n}{Bernoulli}\PYG{p}{(}
        \PYG{l+s+s2}{\PYGZdq{}}\PYG{l+s+s2}{used}\PYG{l+s+s2}{\PYGZdq{}}\PYG{p}{,}
        \PYG{n}{logit\PYGZus{}p}\PYG{o}{=}\PYG{n}{eta}\PYG{p}{,}
        \PYG{n}{observed}\PYG{o}{=}\PYG{n}{y\PYGZus{}data}\PYG{p}{,}
        \PYG{n}{dims}\PYG{o}{=}\PYG{l+s+s2}{\PYGZdq{}}\PYG{l+s+s2}{obs}\PYG{l+s+s2}{\PYGZdq{}}\PYG{p}{,}
    \PYG{p}{)}

    \PYG{n}{inference} \PYG{o}{=} \PYG{n}{pm}\PYG{o}{.}\PYG{n}{sample}\PYG{p}{(}
        \PYG{n}{draws}\PYG{o}{=}\PYG{l+m+mi}{2\PYGZus{}000}\PYG{p}{,}
        \PYG{n}{tune}\PYG{o}{=}\PYG{l+m+mi}{2\PYGZus{}000}\PYG{p}{,}
        \PYG{n}{chains}\PYG{o}{=}\PYG{l+m+mi}{4}\PYG{p}{,}
        \PYG{n}{target\PYGZus{}accept}\PYG{o}{=}\PYG{l+m+mf}{0.95}\PYG{p}{,}
        \PYG{n}{random\PYGZus{}seed}\PYG{o}{=}\PYG{l+m+mi}{42}\PYG{p}{,}
    \PYG{p}{)}
\end{sphinxVerbatim}

\sphinxAtStartPar
The population slope \sphinxcode{\sphinxupquote{beta}} describes the average selection response, while \sphinxcode{\sphinxupquote{sigma\_beta}} quantifies the degree to which that response varies among individuals. Partial pooling regularizes poorly sampled individuals and propagates uncertainty through population\sphinxhyphen{} and individual\sphinxhyphen{}level estimates.

\sphinxAtStartPar
As in the frequentist used\textendash{}available model, the intercept depends on the availability\sphinxhyphen{}sampling design. Individual slopes and their population distribution are generally the principal ecological targets.

\begin{sphinxadmonition}{note}{Note:}
\sphinxAtStartPar
The hierarchical model is currently presented as an optional extension rather than a stabilized hrHSA interface. A future Bayesian module will connect design\sphinxhyphen{}matrix construction, prior specification, posterior diagnostics, spatial prediction, and posterior predictive validation within the package workflow.
\end{sphinxadmonition}

\sphinxstepscope


\chapter{Theory}
\label{\detokenize{theory:theory}}\label{\detokenize{theory::doc}}

\section{1. Foundations}
\label{\detokenize{theory:foundations}}

\subsection{Habitat}
\label{\detokenize{theory:habitat}}
\sphinxAtStartPar
Habitat provides the ecological link between individual movement and population redistribution. By remaining, departing, returning or continuing onwards, animals allocate their time among heterogeneous environments. Repeated over time, these decisions produce patterns of local space use, home\sphinxhyphen{}range formation, migration and seasonal range occupation; aggregated across individuals, they generate population\sphinxhyphen{}level patterns of distribution and density. Habitat selection therefore connects the relocation of the individual with the organisation of the population (Fretwell \& Lucas, 1969; Morris, 2003; Nathan et al., 2008; Mueller \& Fagan, 2008; Matthiopoulos et al., 2015).

\sphinxAtStartPar
Within Habitat Selection Analysis, \sphinxstyleemphasis{habitat} is most usefully understood as a point in environmental space defined by the conjunction of conditions, resources and risks relevant to the focal organism (Morris, 2003; Matthiopoulos et al., 2020; Northrup et al., 2022). Conditions influence physiological or behavioural functioning; resources are required for maintenance, growth or reproduction; and risks reduce expected survival or reproductive success. Habitat is therefore organism\sphinxhyphen{}specific: the same geographic location may present different opportunities and constraints to individuals differing in state, experience or perceptual capacity.

\sphinxAtStartPar
Several related terms must be distinguished. \sphinxstyleemphasis{Use} is the realised allocation of an animal’s time or activity to a habitat. \sphinxstyleemphasis{Occupancy} denotes physical presence within a spatial unit during a defined period. \sphinxstyleemphasis{Availability} describes the habitats that could have been encountered or used, given the spatial, temporal and behavioural constraints acting upon the animal. \sphinxstyleemphasis{Selection} is differential use relative to that availability, whereas \sphinxstyleemphasis{preference}, in its strict sense, describes selection where all alternatives are equally available (Johnson, 1980; Manly et al., 2002; Lele et al., 2013).

\sphinxAtStartPar
For habitat types \(h=1,\ldots,H\), let \(a_h\) denote the proportion of available habitat belonging to type \(h\), and let \(u_h\) denote its proportion of observed use. For \(a_h>0\), the elementary selection ratio is
\begin{equation*}
\begin{split}
r_h = \frac{u_h}{a_h}.
\end{split}
\end{equation*}
\sphinxAtStartPar
Use is proportional to availability when \(r_h=1\), greater than expected when \(r_h>1\), and less than expected when \(r_h<1\). Selection is thus a relationship between use and availability, not an intrinsic property of the habitat. A widespread habitat may contain many recorded locations while receiving less use than expected from its abundance; a rare habitat may contain few observations while nevertheless receiving strongly disproportionate use.


\subsection{The Ideal Free Distribution}
\label{\detokenize{theory:the-ideal-free-distribution}}
\sphinxAtStartPar
Although selection is enacted by individuals, habitat value depends partly upon the population already occupying it. Competitors may deplete resources, interfere with foraging, defend territories or exclude subordinate animals. Population distribution is therefore both a consequence of earlier habitat\sphinxhyphen{}selection decisions and part of the ecological context governing subsequent ones. Let \(N_h\) denote the number or density of animals occupying habitat \(h\), and let \(\Phi_h(N_h)\) denote the expected per\sphinxhyphen{}capita fitness payoff available there. Under competition,
\begin{equation*}
\begin{split}
\frac{\partial \Phi_h(N_h)}{\partial N_h} < 0,
\end{split}
\end{equation*}
\sphinxAtStartPar
so that habitat value declines as density increases. Habitat selection is therefore inherently density dependent (Fretwell \& Lucas, 1969; Rosenzweig, 1981; Morris, 2003; McLoughlin et al., 2010).

\sphinxAtStartPar
The classical formulation of this relationship is the \sphinxstylestrong{Ideal Free Distribution} of Fretwell and Lucas (1969). The term \sphinxstyleemphasis{ideal} assumes that individuals can distinguish the relative payoffs of available habitats; the term \sphinxstyleemphasis{free} assumes that they may move among them without prohibitive costs or exclusion. Consider a population of total size \(N\) distributed among \(H\) habitats, with \(N_h^{*}\) individuals occupying habitat \(h\) at equilibrium:
\begin{equation*}
\begin{split}
\sum_{h=1}^{H} N_h^{*} = N.
\end{split}
\end{equation*}
\sphinxAtStartPar
An ideal\sphinxhyphen{}free equilibrium is reached when no individual could improve its expected fitness by moving elsewhere. All occupied habitats therefore provide the same equilibrium payoff \(\lambda\):
\begin{equation*}
\begin{split}
\Phi_h\left(N_h^{*}\right)
=
\lambda
\qquad
\text{for all } h \text{ with } N_h^{*}>0,
\end{split}
\end{equation*}
\sphinxAtStartPar
while any unoccupied habitat offers no greater return:
\begin{equation*}
\begin{split}
\Phi_h(0)
\leq
\lambda
\qquad
\text{when } N_h^{*}=0.
\end{split}
\end{equation*}
\sphinxAtStartPar
A habitat with a high initial payoff may consequently be occupied first, but the arrival of additional individuals reduces its per\sphinxhyphen{}capita value. Once its payoff falls to that available elsewhere, individuals distribute themselves among habitats in the proportions required to maintain equal expected fitness. Habitats of greater productive capacity may therefore support greater animal densities without providing greater realised fitness at equilibrium: their greater density is precisely what reduces the individual payoff to that available in alternative habitats.

\sphinxAtStartPar
For two habitats, combinations of densities satisfying
\begin{equation*}
\begin{split}
\Phi_1(N_1)=\Phi_2(N_2)
\end{split}
\end{equation*}
\sphinxAtStartPar
form an \sphinxstylestrong{isodar} (Morris, 1987, 1988, 2003). The isodar intercept represents differences in habitat profitability at low density, whereas its slope represents differences in the rate at which fitness declines with crowding. Ideal\sphinxhyphen{}free and isodar theory thus provide a formal bridge between individual habitat choice, density\sphinxhyphen{}dependent competition and population distribution.


\subsection{The Ideal Despotic Distribution}
\label{\detokenize{theory:the-ideal-despotic-distribution}}
\sphinxAtStartPar
The assumptions of the Ideal Free Distribution provide a theoretical reference condition rather than a universal description of nature. Individuals differ in information, state and competitive ability; movement may be costly; and dominant animals may monopolise profitable habitats. Such processes may produce an \sphinxstylestrong{Ideal Despotic Distribution}, in which subordinate individuals occupy less profitable habitats despite the existence of superior alternatives (Fretwell \& Lucas, 1969). Observed distributions may therefore reflect habitat profitability, competition, restricted access and social organisation in combination.

\sphinxAtStartPar
Habitat selection is also hierarchical and scale dependent. Johnson (1980) distinguished selection of a species’ geographical range, placement of an individual home range within that range, use of habitat components within the home range and selection of particular resources within an activity site. A habitat may be selected at one order and avoided at another, or selected for one activity and avoided during another. Availability is therefore an ecological hypothesis concerning the alternatives accessible to an animal at a particular time and scale, rather than a purely geometric property of the study area.

\sphinxAtStartPar
Habitat Selection Analysis begins from these premises. Animals respond to environmental heterogeneity through movement and space use; their decisions produce population distributions; and the resulting densities alter the value of habitats subsequently encountered. The inferential question that follows is how this relationship may be recovered from observations of animal relocations and measured environmental features.


\section{2. Making it real: Towards Inference from Empirical Data}
\label{\detokenize{theory:making-it-real-towards-inference-from-empirical-data}}

\subsection{From Relocations to Selection}
\label{\detokenize{theory:from-relocations-to-selection}}
\sphinxAtStartPar
Habitat selection is a behavioural process, but telemetry records only a finite sample of its spatial consequences. Let \(S_i(t)\) denote the location of individual \(i\) at time \(t\). A tracking device observes this continuous trajectory only at discrete times \(t_{ij}\), producing the relocation series
\begin{equation*}
\begin{split}
\mathcal{D}_i
=
\left\{
\tilde{s}_{ij},t_{ij}
\right\}_{j=1}^{n_i},
\end{split}
\end{equation*}
\sphinxAtStartPar
where \(\tilde{s}_{ij}\) may differ from the true location because of positional error. The observed sample is further conditioned upon the fix schedule and the probability of successfully recording a location.

\sphinxAtStartPar
A concentration of relocations in one part of the landscape may therefore arise because the habitat is selected, because it is readily accessible, because the animal moves slowly within it, or because it repeatedly returns to the same area. Conversely, few relocations may indicate avoidance, restricted access, rapid passage or reduced fix success. Selection cannot be inferred from observations of use alone; the observed distribution must be compared with an explicit model of what was available to the animal.


\subsection{Resource Selection as a Point Process}
\label{\detokenize{theory:resource-selection-as-a-point-process}}
\sphinxAtStartPar
Consider a spatial domain \(\Omega\subset\mathbb{R}^{2}\) containing observed relocations
\begin{equation*}
\begin{split}
\mathcal{S}
=
\left\{
s_1,\ldots,s_n
\right\}.
\end{split}
\end{equation*}
\sphinxAtStartPar
Spatial use may be represented by an \sphinxstylestrong{inhomogeneous Poisson point process}. If \(N(B)\) denotes the number of relocations observed within a region \(B\subseteq\Omega\), then
\begin{equation*}
\begin{split}
N(B)
\sim
\operatorname{Poisson}
\left(
\int_B \lambda(s)\,ds
\right),
\end{split}
\end{equation*}
\sphinxAtStartPar
where \(\lambda(s)>0\) is the spatial intensity of relocations.

\sphinxAtStartPar
Environmental heterogeneity is commonly introduced through a log\sphinxhyphen{}linear intensity:
\begin{equation*}
\begin{split}
\lambda(s)
=
\exp\left[
\beta_0+x(s)\beta
\right],
\end{split}
\end{equation*}
\sphinxAtStartPar
where \(x(s)\) is a row vector of environmental covariates and \(\beta\) is the corresponding vector of coefficients. This may be written as
\begin{equation*}
\begin{split}
\lambda(s)
=
\exp(\beta_0)w\left(x(s)\right),
\end{split}
\end{equation*}
\sphinxAtStartPar
with the selection function
\begin{equation*}
\begin{split}
w(x)
=
\exp(x\beta).
\end{split}
\end{equation*}
\sphinxAtStartPar
The function \(w(x)\) describes how habitat alters the relative intensity of use. It is not the absolute probability that a spatial unit will be occupied or used.

\sphinxAtStartPar
Where locations differ in their accessibility, the intensity may be written more generally as
\begin{equation*}
\begin{split}
\lambda(s)
=
c\,f^{A}(s)w\left(x(s)\right),
\end{split}
\end{equation*}
\sphinxAtStartPar
where \(f^{A}(s)\) is the available distribution and \(c>0\) controls the expected number of observations. Conditional upon observing \(n\) relocations, the corresponding use distribution is
\begin{equation*}
\begin{split}
f^{u}(s)
=
\frac{
f^{A}(s)w\left(x(s)\right)
}{
\displaystyle
\int_{\Omega}
f^{A}(r)w\left(x(r)\right)\,dr
}.
\end{split}
\end{equation*}
\sphinxAtStartPar
The available environment is thus reweighted by the selection function to produce the expected distribution of use.

\sphinxAtStartPar
The intercept has little direct biological meaning in a conventional telemetry analysis because the total number of relocations depends upon the fix schedule, study duration and number of tracked animals. The coefficients in \(\beta\), by contrast, describe relative differences in the intensity of use, conditional upon the defined available distribution.


\subsection{Approximation by Logistic Regression}
\label{\detokenize{theory:approximation-by-logistic-regression}}
\sphinxAtStartPar
The likelihood of an inhomogeneous Poisson point process contains an integral over the full spatial domain:
\begin{equation*}
\begin{split}
L(\beta)
\propto
\exp
\left[
-\int_{\Omega}\lambda(s)\,ds
\right]
\prod_{i=1}^{n}
\lambda(s_i).
\end{split}
\end{equation*}
\sphinxAtStartPar
For realistic landscapes, this integral must usually be approximated numerically. In a use\textendash{}availability design, environmental conditions are evaluated at the observed relocations and at locations sampled from the available distribution.

\sphinxAtStartPar
Let
\begin{equation*}
\begin{split}
Y_j
=
\begin{cases}
1, & \text{if location }j\text{ is used},\\
0, & \text{if location }j\text{ is sampled from availability}.
\end{cases}
\end{split}
\end{equation*}
\sphinxAtStartPar
The combined sample may be fitted using logistic regression:
\begin{equation*}
\begin{split}
\operatorname{logit}
\left[
\Pr(Y_j=1\mid x_j)
\right]
=
\alpha+x_j\beta.
\end{split}
\end{equation*}
\sphinxAtStartPar
Available locations are not biological absences. They represent the environmental distribution from which the animal could have selected and provide a numerical approximation to the point\sphinxhyphen{}process integral. As the available sample becomes sufficiently large, the logistic slopes approach those of the underlying point\sphinxhyphen{}process model (Johnson et al., 2006; Warton \& Shepherd, 2010; Aarts et al., 2012).

\sphinxAtStartPar
The logistic intercept \(\alpha\) depends upon the analyst\sphinxhyphen{}defined ratio of used to available locations and should not ordinarily be interpreted. Predictions should therefore be based upon the exponential selection function
\begin{equation*}
\begin{split}
\hat{w}(x)
=
\exp(x\hat{\beta}),
\end{split}
\end{equation*}
\sphinxAtStartPar
rather than the fitted logistic probability
\begin{equation*}
\begin{split}
\frac{
\exp(\alpha+x\hat{\beta})
}{
1+\exp(\alpha+x\hat{\beta})
}.
\end{split}
\end{equation*}
\sphinxAtStartPar
The number of available locations should be increased until the estimated slopes stabilise. Additional available points improve numerical integration but do not create additional biological replication.


\subsection{Interpretation of an RSF}
\label{\detokenize{theory:interpretation-of-an-rsf}}
\sphinxAtStartPar
For two environmental conditions \(x_1\) and \(x_0\), their relative selection strength is
\begin{equation*}
\begin{split}
\operatorname{RSS}(x_1,x_0)
=
\frac{
\hat{w}(x_1)
}{
\hat{w}(x_0)
}
=
\exp
\left[
(x_1-x_0)\hat{\beta}
\right].
\end{split}
\end{equation*}
\sphinxAtStartPar
For two otherwise identical locations differing by one unit in covariate \(x_j\),
\begin{equation*}
\begin{split}
\operatorname{RSS}
=
\exp(\hat{\beta}_j).
\end{split}
\end{equation*}
\sphinxAtStartPar
Values greater than one indicate greater relative selection for the first condition; values below one indicate lower relative selection. These contrasts remain conditional upon the other model terms, the available domain and the scale of analysis.

\sphinxAtStartPar
Interactions, transformations and nonlinear effects should be interpreted through explicit contrasts rather than through isolated coefficients. A mapped RSF surface,
\begin{equation*}
\begin{split}
\hat{w}\left(x(s)\right)
=
\exp\left[x(s)\hat{\beta}\right],
\end{split}
\end{equation*}
\sphinxAtStartPar
has an arbitrary absolute scale. It may be normalised within a defined available domain to obtain
\begin{equation*}
\begin{split}
\hat{f}^{u}(s)
=
\frac{
f^{A}(s)\hat{w}\left(x(s)\right)
}{
\displaystyle
\int_{\Omega}
f^{A}(r)\hat{w}\left(x(r)\right)\,dr
},
\end{split}
\end{equation*}
\sphinxAtStartPar
but the resulting distribution is specific to that domain.


\subsection{Individuals and Autocorrelation}
\label{\detokenize{theory:individuals-and-autocorrelation}}
\sphinxAtStartPar
Telemetry studies may contain many relocations but comparatively few animals. Locations from the same individual are repeated observations of one behavioural process and should not be treated as independent population replicates.

\sphinxAtStartPar
Individual variation may be represented hierarchically:
\begin{equation*}
\begin{split}
w_i(x)
=
\exp(x\beta_i),
\end{split}
\end{equation*}
\sphinxAtStartPar
with
\begin{equation*}
\begin{split}
\beta_i
\sim
\mathcal{N}(\mu,\Sigma),
\end{split}
\end{equation*}
\sphinxAtStartPar
where \(\mu\) describes the population\sphinxhyphen{}level response and \(\Sigma\) describes variation among individuals. Random slopes are important because animals may differ in the direction or strength of selection. A random intercept alone does not represent such behavioural variation.

\sphinxAtStartPar
Successive relocations are also serially dependent. An animal’s current position constrains its next position, and nearby observations commonly share similar environmental conditions. Ignoring this dependence generally leads to underestimated uncertainty.

\sphinxAtStartPar
Autocorrelation should not automatically be removed by thinning the data until successive points appear independent. Residence, return and slow movement may themselves be biologically meaningful. More defensible approaches include hierarchical models, robust standard errors, block bootstrapping, two\sphinxhyphen{}stage analyses and structured cross\sphinxhyphen{}validation.


\subsection{Model Evaluation}
\label{\detokenize{theory:model-evaluation}}
\sphinxAtStartPar
Evaluation should be based upon held\sphinxhyphen{}out data and should reflect the intended form of prediction. Interpolation within known individuals, prediction to later periods and transfer to new individuals are distinct tasks.

\sphinxAtStartPar
A commonly used measure for presence\sphinxhyphen{}only or use\textendash{}availability predictions is the \sphinxstylestrong{continuous Boyce index}. Let \(\hat{w}(s)\) denote the predicted RSF value in an evaluation dataset. The prediction range is divided into intervals or moving windows \(I_k\). For each interval, calculate the proportion of held\sphinxhyphen{}out used locations,
\begin{equation*}
\begin{split}
P_k
=
\frac{
\text{held-out used locations in }I_k
}{
\text{all held-out used locations}
},
\end{split}
\end{equation*}
\sphinxAtStartPar
and the proportion of available locations,
\begin{equation*}
\begin{split}
E_k
=
\frac{
\text{available locations in }I_k
}{
\text{all available locations}
}.
\end{split}
\end{equation*}
\sphinxAtStartPar
Their ratio is
\begin{equation*}
\begin{split}
R_k
=
\frac{P_k}{E_k}.
\end{split}
\end{equation*}
\sphinxAtStartPar
The Boyce index is the Spearman rank correlation between the representative prediction value \(v_k\) and the predicted\sphinxhyphen{}to\sphinxhyphen{}expected ratio:
\begin{equation*}
\begin{split}
B
=
\operatorname{cor}_{\mathrm{S}}
\left(
v_k,R_k
\right).
\end{split}
\end{equation*}
\sphinxAtStartPar
Values approaching \(1\) indicate that greater predictions correspond to greater use relative to availability. Values near \(0\) indicate little monotonic relationship, while negative values indicate that held\sphinxhyphen{}out locations occur disproportionately in areas assigned low predictions.

\sphinxAtStartPar
The Boyce index evaluates ranking rather than complete probabilistic calibration. It does not preserve the magnitude or geographic location of prediction errors and may depend upon window width, ties and sample size. It should therefore be accompanied by the full predicted\sphinxhyphen{}to\sphinxhyphen{}expected curve, observed and expected use across prediction quantiles, and spatial inspection of held\sphinxhyphen{}out residual patterns.


\subsection{Structured Cross\sphinxhyphen{}Validation}
\label{\detokenize{theory:structured-cross-validation}}
\sphinxAtStartPar
Randomly dividing relocations among folds is generally unsuitable for telemetry data. Neighbouring observations from the same movement bout may enter both the training and evaluation sets, producing information leakage and overly optimistic performance estimates.

\sphinxAtStartPar
For prediction to another period within the same animals, data should be divided into contiguous \sphinxstylestrong{temporal blocks}. Blocks should be sufficiently long to reduce dependence between training and evaluation observations. Where forecasting is the objective, models should be fitted to earlier periods and evaluated on later ones.

\sphinxAtStartPar
For inference to the wider population, \sphinxstylestrong{leave\sphinxhyphen{}one\sphinxhyphen{}individual\sphinxhyphen{}out cross\sphinxhyphen{}validation} is more appropriate. All relocations from one animal are withheld, the model is fitted to the remaining animals, and its predictions are evaluated against the held\sphinxhyphen{}out individual. Repeating this procedure assesses whether the estimated population response transfers to animals not represented during fitting.

\sphinxAtStartPar
Where both temporal transfer and population transfer matter, the two schemes may be combined. An outer leave\sphinxhyphen{}one\sphinxhyphen{}individual\sphinxhyphen{}out loop can assess generalisation to new animals, while temporally blocked validation within the training animals can guide model development.

\sphinxAtStartPar
Validation scores should be reported by individual and fold. A single pooled score may be dominated by animals contributing the greatest number of relocations.


\subsection{From Resource Selection to Step Selection}
\label{\detokenize{theory:from-resource-selection-to-step-selection}}
\sphinxAtStartPar
A conventional RSF usually compares every relocation with a common availability distribution, such as an individual’s home range or the study area. At fine temporal scales, this assumption becomes difficult to defend. A location may lie within the home range but remain unreachable during the interval between two fixes.

\sphinxAtStartPar
Movement therefore determines availability. The locations accessible at time \(t+1\) depend upon the animal’s position at time \(t\), the elapsed time, its movement capacity and its previous direction of travel.

\sphinxAtStartPar
A \sphinxstylestrong{step} is the displacement between consecutive relocations:
\begin{equation*}
\begin{split}
s_t
\longrightarrow
s_{t+1}.
\end{split}
\end{equation*}
\sphinxAtStartPar
Its length is
\begin{equation*}
\begin{split}
l_t
=
\left\|
s_{t+1}-s_t
\right\|,
\end{split}
\end{equation*}
\sphinxAtStartPar
and its turning angle \(\theta_t\) is the change in bearing relative to the preceding step.

\sphinxAtStartPar
In a \sphinxstylestrong{Step Selection Function}, each observed step is paired with available steps beginning at the same location. Their lengths and turning angles are conventionally sampled from the empirical distributions of observed movements. Habitat conditions at the observed endpoint or along the observed path are then compared with those of the available alternatives.

\sphinxAtStartPar
Let
\begin{equation*}
\begin{split}
\mathcal{C}_t
=
\left\{
s_{t+1}^{(0)},
s_{t+1}^{(1)},
\ldots,
s_{t+1}^{(K)}
\right\}
\end{split}
\end{equation*}
\sphinxAtStartPar
denote the choice set for step \(t\), where \(s_{t+1}^{(0)}\) is the observed endpoint. Under conditional logistic regression,
\begin{equation*}
\begin{split}
\Pr
\left(
j\mid\mathcal{C}_t
\right)
=
\frac{
\exp\left(z_{tj}\gamma\right)
}{
\displaystyle
\sum_{k=0}^{K}
\exp\left(z_{tk}\gamma\right)
}.
\end{split}
\end{equation*}
\sphinxAtStartPar
An SSF therefore asks which characteristics distinguished the chosen step from the alternatives reachable from the same starting point.


\subsection{Integrated Step Selection}
\label{\detokenize{theory:integrated-step-selection}}
\sphinxAtStartPar
Empirical step\sphinxhyphen{}length and turning\sphinxhyphen{}angle distributions have already been shaped by habitat selection. Short observed steps, for example, may reflect intrinsically slow movement or prolonged residence in selected habitat. Sampling available movements directly from the observed distributions therefore risks confounding movement with selection.

\sphinxAtStartPar
An \sphinxstylestrong{integrated Step Selection Function} estimates the movement and selection processes jointly. Let \(\phi\) denote the baseline movement kernel and \(w\) the habitat\sphinxhyphen{}selection function. The transition density is
\begin{equation*}
\begin{split}
p
\left(
s_{t+1}
\mid
s_t,s_{t-1}
\right)
=
\frac{
\phi
\left(
s_{t+1}
\mid
s_t,s_{t-1};\theta
\right)
w
\left(
x(s_{t+1});\beta
\right)
}{
\displaystyle
\int_{\Omega}
\phi
\left(
r
\mid
s_t,s_{t-1};\theta
\right)
w
\left(
x(r);\beta
\right)\,dr
}.
\end{split}
\end{equation*}
\sphinxAtStartPar
The movement kernel determines which locations can be reached and with what baseline probability. The selection function then reweights those locations according to habitat. Movement and selection together determine the next relocation.

\sphinxAtStartPar
Terms such as \(l\), \(\log l\) and \(\cos\theta\) allow parameters of step\sphinxhyphen{}length and turning\sphinxhyphen{}angle distributions to be estimated alongside habitat\sphinxhyphen{}selection coefficients. The resulting model separates the baseline movement process from the effects of environmental conditions more explicitly than a conventional SSF (Avgar et al., 2016).


\subsection{Interpretation Through Simulation}
\label{\detokenize{theory:interpretation-through-simulation}}
\sphinxAtStartPar
An SSF or iSSF coefficient describes a local contrast among alternatives available from the same starting point. For two candidates \(a\) and \(b\),
\begin{equation*}
\begin{split}
\frac{
\Pr(a\mid\mathcal{C}_t)
}{
\Pr(b\mid\mathcal{C}_t)
}
=
\exp
\left[
\eta_{ta}-\eta_{tb}
\right].
\end{split}
\end{equation*}
\sphinxAtStartPar
Such contrasts remain directly interpretable. The long\sphinxhyphen{}term consequences of the coefficients are less immediate because each selected endpoint determines the next choice set. Movement and selection therefore alter both the present choice and the locations available in the future.

\sphinxAtStartPar
Predicted utilisation distributions, residence times, crossing probabilities and home\sphinxhyphen{}range geometry should consequently be obtained through simulation of the complete transition process. At each step, candidate movements are drawn from the fitted movement kernel, weighted according to habitat, and sampled to produce the next location. Repetition generates trajectories whose emergent spatial patterns can be compared with held\sphinxhyphen{}out observations.

\sphinxAtStartPar
RSFs, SSFs and iSSFs therefore represent a progression in the treatment of availability. An RSF relates use to a broader and approximately static available distribution. An SSF conditions availability upon the animal’s previous location. An iSSF jointly estimates the movement process generating those local alternatives and the habitat\sphinxhyphen{}selection process distinguishing among them. The appropriate framework depends upon the temporal scale of the data and the inferential question being addressed.


\chapter{Towards Bayesian Inference}
\label{\detokenize{theory:towards-bayesian-inference}}
\sphinxAtStartPar
The preceding models commonly proceed as though recorded locations were exact, environmental covariates were known without error, individuals differed only through random sampling, and habitat\sphinxhyphen{}selection relationships remained constant through time. Each assumption is convenient; none is generally true. A Bayesian analysis is not the only means of relaxing them, but it provides a coherent framework in which uncertain observations, latent ecological processes, individual heterogeneity and temporal change may be represented jointly and propagated into the final inference.

\sphinxAtStartPar
Let \(\mathcal{D}\) denote the observed data, \(\mathcal{Z}\) the unobserved ecological quantities from which those data arose, and \(\Theta\) the model parameters. Bayesian inference proceeds from the posterior distribution
\begin{equation*}
\begin{split}
p(\mathcal{Z},\Theta\mid\mathcal{D})
\propto
p(\mathcal{D}\mid\mathcal{Z},\Theta)
p(\mathcal{Z}\mid\Theta)
p(\Theta).
\end{split}
\end{equation*}
\sphinxAtStartPar
The first term describes the observation process, the second the ecological process and the third the prior information assigned to its parameters. Rather than replacing uncertain quantities by single estimates, the posterior integrates over their plausible values. This distinction is central to three objectives of \sphinxstylestrong{hrHSA}: propagating uncertainty, estimating individual variation and detecting change.


\section{Propagating Locational and Environmental Uncertainty}
\label{\detokenize{theory:propagating-locational-and-environmental-uncertainty}}
\sphinxAtStartPar
Telemetry records an estimate of an animal’s location rather than the location itself. Let \(s_{ij}\) denote the true location of individual \(i\) at observation \(j\), and let \(\tilde{s}_{ij}\) denote the recorded GNSS fix. A simple observation model is
\begin{equation*}
\begin{split}
\tilde{s}_{ij}
\mid
s_{ij},\Sigma_{ij}
\sim
\mathcal{N}_2
\left(
s_{ij},
\Sigma_{ij}
\right),
\end{split}
\end{equation*}
\sphinxAtStartPar
where \(\Sigma_{ij}\) describes the magnitude and orientation of locational uncertainty. Other distributions may be substituted where errors are heavy\sphinxhyphen{}tailed, device\sphinxhyphen{}specific or otherwise non\sphinxhyphen{}Gaussian. The ecological model is then evaluated at the latent location \(s_{ij}\) rather than treating \(\tilde{s}_{ij}\) as exact (Jonsen et al., 2005; Patterson et al., 2008; Hooten et al., 2017).

\sphinxAtStartPar
Locational inaccuracy must be distinguished from failed fixes. A recorded fix may be spatially imprecise, whereas a failed fix produces no location at all. If fix success depends upon canopy cover, topography or other habitat features, the observed relocations constitute a biased sample of the underlying path. The observation process may therefore require both a model for positional error and a model for the probability of detection or successful acquisition (Frair et al., 2004; Nielson et al., 2009).

\sphinxAtStartPar
Environmental data are likewise imperfect. Let \(x^{*}(s,t)\) denote the environmental field relevant to the animal and let \(\tilde{x}_g\) denote the value represented by raster cell \(g\), covering area \(A_g\). A continuous raster variable may be viewed as an imperfect spatial aggregate:
\begin{equation*}
\begin{split}
\tilde{x}_g
=
\frac{1}{|A_g|}
\int_{A_g}
x^{*}(s,t)\,ds
+
\epsilon_g,
\end{split}
\end{equation*}
\sphinxAtStartPar
where \(\epsilon_g\) represents measurement, interpolation or classification error. A raster cell therefore describes an average or assigned class over an area, while the animal occupies a location within it. This difference in spatial support becomes consequential where the environmental field varies substantially within cells or where locational uncertainty is large relative to raster resolution.

\sphinxAtStartPar
Raster uncertainty arises from several sources: sensor and classification error, temporal mismatch, interpolation, reprojection, resampling and the averaging of heterogeneous conditions within pixels. Grain size is not merely another random error term. It determines the spatial support over which habitat is represented and may therefore alter the ecological relationship being estimated. Uncertainty about grain should be addressed through biologically motivated multiscale models or sensitivity analyses rather than concealed within a single residual variance (Northrup et al., 2022).

\sphinxAtStartPar
A hierarchical model may combine these uncertainties schematically as
\begin{equation*}
\begin{split}
\begin{aligned}
\tilde{s}_{ij}
&\sim
p\left(
\tilde{s}_{ij}\mid s_{ij},\psi_s
\right),\\
\tilde{x}
&\sim
p\left(
\tilde{x}\mid x^{*},\psi_x
\right),\\
y_{ij}
&\sim
p\left(
y_{ij}\mid s_{ij},x^{*},\beta_i
\right),
\end{aligned}
\end{split}
\end{equation*}
\sphinxAtStartPar
where \(\psi_s\) and \(\psi_x\) govern locational and environmental uncertainty. The posterior distribution of \(\beta_i\) then reflects uncertainty not only in sampling, but also in the positions and environmental conditions upon which inference depends.

\sphinxAtStartPar
This propagation is especially important for derived quantities. Selection surfaces, utilisation distributions and simulated trajectories should not be calculated solely from posterior mean coefficients and a single environmental raster. They may instead be generated repeatedly from posterior draws of locations, environmental fields and model parameters, producing a distribution of predictions rather than one deceptively exact map.


\section{Individual Variation as Biological Information}
\label{\detokenize{theory:individual-variation-as-biological-information}}
\sphinxAtStartPar
A population\sphinxhyphen{}level coefficient describes a central tendency; it does not imply that all individuals respond alike. Animals may differ in habitat selection because of sex, age, reproductive state, experience, physiology, social status, competitive ability or behavioural phenotype. Pooling these differences into a single coefficient may obscure opposing responses and understate uncertainty at the population level (Leclerc et al., 2016; Muff et al., 2020; Northrup et al., 2022).

\sphinxAtStartPar
Individual\sphinxhyphen{}specific selection coefficients may be represented hierarchically as
\begin{equation*}
\begin{split}
\beta_i
=
\mu_{\beta}
+
Z_i\gamma
+
b_i,
\end{split}
\end{equation*}
\sphinxAtStartPar
with
\begin{equation*}
\begin{split}
b_i
\sim
\mathcal{N}
\left(
0,\Sigma_{\beta}
\right).
\end{split}
\end{equation*}
\sphinxAtStartPar
Here, \(\mu_{\beta}\) is the population\sphinxhyphen{}level mean response, \(Z_i\gamma\) describes systematic differences associated with measured individual attributes, and \(b_i\) represents residual variation among individuals. The covariance matrix \(\Sigma_{\beta}\) quantifies both the magnitude of individual variation and correlations among selection responses.

\sphinxAtStartPar
This model partially pools information across animals. Individuals with abundant and informative data retain strongly individual estimates, whereas poorly sampled individuals are drawn towards the population distribution. Partial pooling avoids the two extremes of treating every relocation as an independent population replicate and fitting entirely separate models whose imprecise estimates are subsequently treated as exact.

\sphinxAtStartPar
Individual variation also changes population\sphinxhyphen{}level prediction. Because the selection function is nonlinear,
\begin{equation*}
\begin{split}
\operatorname{E}_i
\left[
\exp(x\beta_i)
\right]
\neq
\exp
\left[
x\operatorname{E}_i(\beta_i)
\right].
\end{split}
\end{equation*}
\sphinxAtStartPar
The prediction obtained from the mean individual is therefore not generally equal to the mean prediction across individuals. Population\sphinxhyphen{}level maps should integrate over the estimated distribution of individual responses rather than merely substitute \(\mu_{\beta}\) into the selection function.

\sphinxAtStartPar
Among\sphinxhyphen{}individual variation is not statistical debris around a population mean. It is a property of the population and may reveal individual specialisation, alternative behavioural tactics or differences in environmental sensitivity. Phenotypic variation is also the material upon which natural selection can act. Variation in estimated habitat\sphinxhyphen{}selection behaviour is not, by itself, evidence of adaptive potential: an evolutionary response additionally requires that differences are sufficiently persistent, heritable and associated with fitness. Hierarchical habitat\sphinxhyphen{}selection models can establish and quantify behavioural variation, while pedigrees, genomic information, repeated observations and demographic data are required to distinguish its genetic, environmental and fitness\sphinxhyphen{}related components (Leclerc et al., 2016; Gervais et al., 2022).


\section{Detecting Change}
\label{\detokenize{theory:detecting-change}}
\sphinxAtStartPar
Conventional selection models assume that the coefficient vector \(\beta\) remains constant over the period of inference. Yet selection may vary with season, life\sphinxhyphen{}history stage, population density, resource availability, disturbance, learning or climatic conditions. In a changing environment, it is therefore insufficient to ask only which habitats are selected on average. We must also ask whether, when and among whom the relationship is changing.

\sphinxAtStartPar
A time\sphinxhyphen{}varying selection function may be written as
\begin{equation*}
\begin{split}
w_{i,t}(x)
=
\exp
\left(
x\beta_{i,t}
\right).
\end{split}
\end{equation*}
\sphinxAtStartPar
Where change is expected to be gradual, the coefficients may evolve according to a stochastic process such as
\begin{equation*}
\begin{split}
\beta_{i,t}
\sim
\mathcal{N}
\left(
\beta_{i,t-1},
Q
\right),
\end{split}
\end{equation*}
\sphinxAtStartPar
where \(Q\) controls the rate and covariance of temporal change. Smaller values of \(Q\) imply comparatively stable selection, whereas larger values permit more rapid variation. Alternative models may impose seasonal periodicity, smooth temporal trends or dependence upon measured environmental drivers.

\sphinxAtStartPar
Where an abrupt transition is expected, a change\sphinxhyphen{}point model may instead be used:
\begin{equation*}
\begin{split}
\beta_{i,t}
=
\begin{cases}
\beta_i^{(1)}, & t < \tau_i,\\
\beta_i^{(2)}, & t \geq \tau_i,
\end{cases}
\end{split}
\end{equation*}
\sphinxAtStartPar
where \(\tau_i\) is the unknown time of change. Bayesian inference yields a posterior distribution for \(\tau_i\) rather than forcing the transition to coincide with an analyst\sphinxhyphen{}defined season or calendar date.

\sphinxAtStartPar
Temporal change and individual variation may be separated through
\begin{equation*}
\begin{split}
\beta_{i,t}
=
\mu_t
+
b_i
+
u_{i,t},
\end{split}
\end{equation*}
\sphinxAtStartPar
where \(\mu_t\) is the population\sphinxhyphen{}wide trajectory, \(b_i\) is a persistent individual deviation and \(u_{i,t}\) describes within\sphinxhyphen{}individual change. This distinction separates population\sphinxhyphen{}wide redistribution from differences in individual timing and response. It also prevents changes in the composition of the sampled population from being mistaken for behavioural change within individuals.

\sphinxAtStartPar
Posterior distributions permit direct statements concerning change. For a coefficient \(\beta_j\), one may calculate
\begin{equation*}
\begin{split}
\Pr
\left(
\beta_{j,t_2}
-
\beta_{j,t_1}
>
0
\mid
\mathcal{D}
\right),
\end{split}
\end{equation*}
\sphinxAtStartPar
or the posterior probability that a change exceeds a biologically meaningful threshold. The result describes the magnitude, direction and uncertainty of change rather than reducing inference to a sequence of separate significance tests.

\sphinxAtStartPar
This capacity is particularly relevant under global change. Changes in climate, land use and human disturbance may alter both the distribution of available habitat and the manner in which animals respond to it. A dynamic model can distinguish a changing environmental landscape from a changing selection relationship, provided that both used and available conditions are represented through time. It may reveal gradual adjustment, abrupt displacement, increased variability or divergent responses among individuals.

\sphinxAtStartPar
Temporal change in a coefficient does not, however, identify its cause. Apparent non\sphinxhyphen{}stationarity may arise from season, reproduction, ageing, population density, behavioural state, changing availability or alteration of the observation process. Attribution to global change requires explicit environmental hypotheses, suitable temporal replication and, where possible, comparison across populations or landscapes. Dynamic Bayesian models expose change and quantify its uncertainty; they do not replace ecological reasoning.


\section{Posterior Prediction and Model Evaluation}
\label{\detokenize{theory:posterior-prediction-and-model-evaluation}}
\sphinxAtStartPar
The principal practical advantage of the Bayesian formulation is that uncertainty can be carried into every derived prediction. Let \(\Theta\) collect the parameters and latent states of the fitted model. The posterior predictive distribution is
\begin{equation*}
\begin{split}
p
\left(
\tilde{\mathcal{D}}
\mid
\mathcal{D}
\right)
=
\int
p
\left(
\tilde{\mathcal{D}}
\mid
\Theta
\right)
p
\left(
\Theta
\mid
\mathcal{D}
\right)
\,d\Theta.
\end{split}
\end{equation*}
\sphinxAtStartPar
Each posterior draw represents one plausible state of the ecological system. Repeated prediction or simulation therefore propagates uncertainty in locations, environmental covariates, individual responses and temporal dynamics into selection maps, utilisation distributions and simulated movement paths.

\sphinxAtStartPar
Bayesian inference does not remove the need for validation. Priors should be examined through prior\sphinxhyphen{}predictive simulation, fitted models through posterior\sphinxhyphen{}predictive checks, and predictive performance through the temporally blocked and leave\sphinxhyphen{}one\sphinxhyphen{}individual\sphinxhyphen{}out procedures described in the preceding chapter. Weak identifiability, an inappropriate availability distribution or an inadequate ecological model cannot be repaired by increasingly elaborate computation.

\sphinxAtStartPar
The purpose of the Bayesian extension is therefore not complexity for its own sake. It is to bring the inferential model into closer correspondence with the ecological system: locations and environments are observed imperfectly, populations consist of heterogeneous individuals, and selection may change through time. Representing these processes jointly permits uncertainty to remain visible, individual variation to become an object of inference and ecological change to be detected rather than averaged away.


\section{Principal References}
\label{\detokenize{theory:principal-references}}
\sphinxAtStartPar
Auger\sphinxhyphen{}Méthé, M., Newman, K., Cole, D., Empacher, F., Gryba, R., King, A. A., Leos\sphinxhyphen{}Barajas, V., Mills Flemming, J., Nielsen, A., Petris, G. \& Thomas, L. (2021). A guide to state\sphinxhyphen{}space modeling of ecological time series. \sphinxstyleemphasis{Ecological Monographs}, \sphinxstylestrong{91}, e01470.

\sphinxAtStartPar
Dejeante, R., Valeix, M. \& Chamaillé\sphinxhyphen{}Jammes, S. (2024). Time\sphinxhyphen{}varying habitat selection analysis: a model and applications for studying diel, seasonal, and post\sphinxhyphen{}release changes. \sphinxstyleemphasis{Ecology}, \sphinxstylestrong{105}, e4233.

\sphinxAtStartPar
Frair, J. L., Nielsen, S. E., Merrill, E. H., Lele, S. R., Boyce, M. S., Munro, R. H. M., Stenhouse, G. B. \& Beyer, H. L. (2004). Removing GPS collar bias in habitat selection studies. \sphinxstyleemphasis{Journal of Applied Ecology}, \sphinxstylestrong{41}, 201\textendash{}212.

\sphinxAtStartPar
Gervais, L., Morellet, N., David, I., Hewison, M., Réale, D., Goulard, M., Chaval, Y., Lourtet, B., Cargnelutti, B., Merlet, J., Quéméré, E. \& Pujol, B. (2022). Quantifying heritability and estimating evolutionary potential in the wild when individuals that share genes also share environments. \sphinxstyleemphasis{Journal of Animal Ecology}, \sphinxstylestrong{91}, 1239\textendash{}1250.

\sphinxAtStartPar
Hooten, M. B., Johnson, D. S., McClintock, B. T. \& Morales, J. M. (2017). \sphinxstyleemphasis{Animal Movement: Statistical Models for Telemetry Data}. Boca Raton: CRC Press.

\sphinxAtStartPar
Jonsen, I. D., Flemming, J. M. \& Myers, R. A. (2005). Robust state\sphinxhyphen{}space modeling of animal movement data. \sphinxstyleemphasis{Ecology}, \sphinxstylestrong{86}, 2874\textendash{}2880.

\sphinxAtStartPar
Leclerc, M., Vander Wal, E., Zedrosser, A., Swenson, J. E., Kindberg, J. \& Pelletier, F. (2016). Quantifying consistent individual differences in habitat selection. \sphinxstyleemphasis{Oecologia}, \sphinxstylestrong{180}, 697\textendash{}705.

\sphinxAtStartPar
Muff, S., Signer, J. \& Fieberg, J. (2020). Accounting for individual\sphinxhyphen{}specific variation in habitat\sphinxhyphen{}selection studies: efficient estimation of mixed\sphinxhyphen{}effects models using Bayesian or frequentist computation. \sphinxstyleemphasis{Journal of Animal Ecology}, \sphinxstylestrong{89}, 80\textendash{}92.

\sphinxAtStartPar
Nielson, R. M., Manly, B. F. J., McDonald, L. L., Sawyer, H. \& McDonald, T. L. (2009). Estimating habitat selection when GPS fix success is less than 100\%. \sphinxstyleemphasis{Ecology}, \sphinxstylestrong{90}, 2956\textendash{}2962.

\sphinxAtStartPar
Northrup, J. M., Vander Wal, E., Bonar, M., Fieberg, J., Laforge, M. P., Leclerc, M., Prokopenko, C. M. \& Gerber, B. D. (2022). Conceptual and methodological advances in habitat\sphinxhyphen{}selection modeling: guidelines for ecology and evolution. \sphinxstyleemphasis{Ecological Applications}, \sphinxstylestrong{32}, e02470.

\sphinxAtStartPar
Patterson, T. A., Thomas, L., Wilcox, C., Ovaskainen, O. \& Matthiopoulos, J. (2008). State\sphinxhyphen{}space models of individual animal movement. \sphinxstyleemphasis{Trends in Ecology \& Evolution}, \sphinxstylestrong{23}, 87\textendash{}94.

\sphinxAtStartPar
Aarts, G., Fieberg, J. \& Matthiopoulos, J. (2012). Comparative interpretation of count, presence\textendash{}absence and point methods for species distribution models. \sphinxstyleemphasis{Methods in Ecology and Evolution}, \sphinxstylestrong{3}, 177\textendash{}187.

\sphinxAtStartPar
Avgar, T., Potts, J. R., Lewis, M. A. \& Boyce, M. S. (2016). Integrated step selection analysis: bridging the gap between resource selection and animal movement. \sphinxstyleemphasis{Methods in Ecology and Evolution}, \sphinxstylestrong{7}, 619\textendash{}630.

\sphinxAtStartPar
Boyce, M. S., Vernier, P. R., Nielsen, S. E. \& Schmiegelow, F. K. A. (2002). Evaluating Resource Selection Functions. \sphinxstyleemphasis{Ecological Modelling}, \sphinxstylestrong{157}, 281\textendash{}300.

\sphinxAtStartPar
Fieberg, J., Matthiopoulos, J., Hebblewhite, M., Boyce, M. S. \& Frair, J. L. (2010). Correlation and studies of habitat selection: problem, red herring or opportunity? \sphinxstyleemphasis{Philosophical Transactions of the Royal Society B}, \sphinxstylestrong{365}, 2233\textendash{}2244.

\sphinxAtStartPar
Fieberg, J., Signer, J., Smith, B. \& Avgar, T. (2021). A “how to” guide for interpreting parameters in habitat\sphinxhyphen{}selection analyses. \sphinxstyleemphasis{Journal of Animal Ecology}, \sphinxstylestrong{90}, 1027\textendash{}1043.

\sphinxAtStartPar
Forester, J. D., Im, H. K. \& Rathouz, P. J. (2009). Accounting for animal movement in estimation of Resource Selection Functions: sampling and data analysis. \sphinxstyleemphasis{Ecology}, \sphinxstylestrong{90}, 3554\textendash{}3565.

\sphinxAtStartPar
Fortin, D., Beyer, H. L., Boyce, M. S., Smith, D. W., Duchesne, T. \& Mao, J. S. (2005). Wolves influence elk movements: behavior shapes a trophic cascade in Yellowstone National Park. \sphinxstyleemphasis{Ecology}, \sphinxstylestrong{86}, 1320\textendash{}1330.

\sphinxAtStartPar
Hirzel, A. H., Le Lay, G., Helfer, V., Randin, C. \& Guisan, A. (2006). Evaluating the ability of habitat\sphinxhyphen{}suitability models to predict species presences. \sphinxstyleemphasis{Ecological Modelling}, \sphinxstylestrong{199}, 142\textendash{}152.

\sphinxAtStartPar
Johnson, C. J., Nielsen, S. E., Merrill, E. H., McDonald, T. L. \& Boyce, M. S. (2006). Resource Selection Functions based on use\textendash{}availability data: theoretical motivation and evaluation methods. \sphinxstyleemphasis{Journal of Wildlife Management}, \sphinxstylestrong{70}, 347\textendash{}357.

\sphinxAtStartPar
Northrup, J. M., Hooten, M. B., Anderson, C. R. \& Wittemyer, G. (2013). Practical guidance on characterizing availability in Resource Selection Functions under a use\textendash{}availability design. \sphinxstyleemphasis{Ecology}, \sphinxstylestrong{94}, 1456\textendash{}1463.

\sphinxAtStartPar
Roberts, D. R. et al. (2017). Cross\sphinxhyphen{}validation strategies for data with temporal, spatial, hierarchical, or phylogenetic structure. \sphinxstyleemphasis{Ecography}, \sphinxstylestrong{40}, 913\textendash{}929.

\sphinxAtStartPar
Warton, D. I. \& Shepherd, L. C. (2010). Poisson point\sphinxhyphen{}process models solve the “pseudo\sphinxhyphen{}absence problem” for presence\sphinxhyphen{}only data in ecology. \sphinxstyleemphasis{The Annals of Applied Statistics}, \sphinxstylestrong{4}, 1383\textendash{}1402.

\sphinxstepscope


\chapter{Implementation}
\label{\detokenize{implementation:implementation}}\label{\detokenize{implementation::doc}}

\section{Design Principles}
\label{\detokenize{implementation:design-principles}}
\sphinxAtStartPar
\sphinxstylestrong{hrHSA} is implemented as a scientific software package rather than as a collection of study\sphinxhyphen{}specific scripts. Its architecture is guided by four principles:
\begin{enumerate}
\sphinxsetlistlabels{\arabic}{enumi}{enumii}{}{.}%
\item {} 
\sphinxAtStartPar
\sphinxstylestrong{Reproducibility} — every transformation between the original data and the final inference should be explicit and recoverable.

\item {} 
\sphinxAtStartPar
\sphinxstylestrong{Modularity} — data preparation, feature construction, statistical inference, validation and prediction should remain separable.

\item {} 
\sphinxAtStartPar
\sphinxstylestrong{Efficiency} — computational resources should be used deliberately, with particular attention to memory access, spatial chunking and repeated raster operations.

\item {} 
\sphinxAtStartPar
\sphinxstylestrong{Scalability} — the same analytical workflow should run on a notebook, a multicore workstation or a SLURM\sphinxhyphen{}managed computing cluster.

\end{enumerate}

\sphinxAtStartPar
The resulting workflow separates scientific decisions from computational execution:

\begin{sphinxVerbatim}[commandchars=\\\{\}]
telemetry and environmental data
        ↓
spatial standardisation
        ↓
definition and sampling of availability
        ↓
environmental feature extraction
        ↓
model\PYGZhy{}matrix construction
        ↓
statistical inference
        ↓
validation and diagnostics
        ↓
spatial prediction or movement simulation
\end{sphinxVerbatim}

\sphinxAtStartPar
Each stage exposes a limited interface and may be inspected, tested or replaced independently.


\section{Package Architecture}
\label{\detokenize{implementation:package-architecture}}
\sphinxAtStartPar
The source code follows the conventional Python \sphinxcode{\sphinxupquote{src}} layout and is divided according to analytical responsibility:

\begin{sphinxVerbatim}[commandchars=\\\{\}]
src/hsa/
├── sampling.py          \PYGZsh{} availability domains and raster sampling
├── features.py          \PYGZsh{} reproducible construction of model matrices
├── types.py             \PYGZsh{} shared specifications and data containers
├── rsf/                 \PYGZsh{} fitting, prediction, validation and cross\PYGZhy{}validation
├── movement/            \PYGZsh{} trajectories, steps, angles and movement kernels
├── remote\PYGZus{}sensing/      \PYGZsh{} optional Earth Observation interfaces
├── compute/             \PYGZsh{} chunking, scalable I/O and distributed execution
└── simulation/          \PYGZsh{} forward movement simulation
\end{sphinxVerbatim}

\sphinxAtStartPar
The package itself remains independent of particular species, reserves, local file systems or Earth Engine projects. Study\sphinxhyphen{}specific assumptions belong in configuration files, examples or downstream analyses. This distinction allows empirical workflows to remain fully documented without embedding them in the reusable statistical core.

\sphinxAtStartPar
Optional functionality is separated into dependency groups. A minimal installation provides the numerical and geospatial foundations of the package, while Earth Engine, distributed computation and HPC support may be installed only where required:

\begin{sphinxVerbatim}[commandchars=\\\{\}]
pip\PYG{+w}{ }install\PYG{+w}{ }\PYGZhy{}e\PYG{+w}{ }.
pip\PYG{+w}{ }install\PYG{+w}{ }\PYGZhy{}e\PYG{+w}{ }\PYG{l+s+s2}{\PYGZdq{}.[earthengine]\PYGZdq{}}
pip\PYG{+w}{ }install\PYG{+w}{ }\PYGZhy{}e\PYG{+w}{ }\PYG{l+s+s2}{\PYGZdq{}.[hpc,io]\PYGZdq{}}
\end{sphinxVerbatim}

\sphinxAtStartPar
Optional libraries are imported lazily inside the functions that require them. Importing \sphinxstylestrong{hrHSA} therefore neither initialises external services nor imposes a distributed\sphinxhyphen{}computing environment upon smaller analyses.


\section{Reproducible Model Construction}
\label{\detokenize{implementation:reproducible-model-construction}}
\sphinxAtStartPar
A fitted model is defined not only by its coefficients, but also by the transformations applied before estimation. \sphinxstylestrong{hrHSA} represents these transformations through a declarative \sphinxcode{\sphinxupquote{FeatureSpec}}:

\begin{sphinxVerbatim}[commandchars=\\\{\}]
\PYG{k+kn}{from}\PYG{+w}{ }\PYG{n+nn}{hsa}\PYG{+w}{ }\PYG{k+kn}{import} \PYG{n}{FeatureSpec}

\PYG{n}{spec} \PYG{o}{=} \PYG{n}{FeatureSpec}\PYG{p}{(}
    \PYG{n}{linear}\PYG{o}{=}\PYG{p}{[}\PYG{l+s+s2}{\PYGZdq{}}\PYG{l+s+s2}{ndvi}\PYG{l+s+s2}{\PYGZdq{}}\PYG{p}{,} \PYG{l+s+s2}{\PYGZdq{}}\PYG{l+s+s2}{distance\PYGZus{}to\PYGZus{}water}\PYG{l+s+s2}{\PYGZdq{}}\PYG{p}{]}\PYG{p}{,}
    \PYG{n}{quadratic}\PYG{o}{=}\PYG{p}{[}\PYG{l+s+s2}{\PYGZdq{}}\PYG{l+s+s2}{distance\PYGZus{}to\PYGZus{}water}\PYG{l+s+s2}{\PYGZdq{}}\PYG{p}{]}\PYG{p}{,}
    \PYG{n}{interactions}\PYG{o}{=}\PYG{p}{[}\PYG{p}{(}\PYG{l+s+s2}{\PYGZdq{}}\PYG{l+s+s2}{ndvi}\PYG{l+s+s2}{\PYGZdq{}}\PYG{p}{,} \PYG{l+s+s2}{\PYGZdq{}}\PYG{l+s+s2}{distance\PYGZus{}to\PYGZus{}water}\PYG{l+s+s2}{\PYGZdq{}}\PYG{p}{)}\PYG{p}{]}\PYG{p}{,}
    \PYG{n}{categorical}\PYG{o}{=}\PYG{p}{[}\PYG{l+s+s2}{\PYGZdq{}}\PYG{l+s+s2}{landcover}\PYG{l+s+s2}{\PYGZdq{}}\PYG{p}{]}\PYG{p}{,}
    \PYG{n}{add\PYGZus{}const}\PYG{o}{=}\PYG{k+kc}{True}\PYG{p}{,}
\PYG{p}{)}
\end{sphinxVerbatim}

\sphinxAtStartPar
Continuous predictors are standardised during fitting. Quadratic and interaction terms are then constructed on the standardised scale. Categorical levels, reference categories and final column order are stored as model metadata.

\begin{sphinxVerbatim}[commandchars=\\\{\}]
\PYG{n}{model}\PYG{p}{,} \PYG{n}{scaler}\PYG{p}{,} \PYG{n}{spec}\PYG{p}{,} \PYG{n}{meta} \PYG{o}{=} \PYG{n}{fit\PYGZus{}rsf}\PYG{p}{(}\PYG{n}{samples}\PYG{p}{,} \PYG{n}{spec}\PYG{p}{)}
\end{sphinxVerbatim}

\sphinxAtStartPar
The model, fitted scaler, feature specification and metadata jointly constitute the fitted analysis. Passing the same objects to validation and prediction ensures that raster surfaces and held\sphinxhyphen{}out observations receive precisely the transformations used during model fitting.

\sphinxAtStartPar
Randomised operations accept explicit seeds. Coordinate reference systems are checked before spatial operations, and functions involving distances require projected coordinates with meaningful linear units. Invalid inputs are rejected early rather than being allowed to produce numerically plausible but ecologically incorrect results.


\section{Geospatial Data and Scalable I/O}
\label{\detokenize{implementation:geospatial-data-and-scalable-i-o}}
\sphinxAtStartPar
Telemetry locations are represented as \sphinxcode{\sphinxupquote{GeoDataFrame}} objects, while environmental fields are stored as labelled \sphinxcode{\sphinxupquote{xarray}} arrays. A conventional environmental stack has the dimensions

\begin{sphinxVerbatim}[commandchars=\\\{\}]
band × y × x
\end{sphinxVerbatim}

\sphinxAtStartPar
and retains its coordinates, projection, spatial resolution and band names throughout the workflow.

\sphinxAtStartPar
Large raster stacks may be stored as \sphinxstylestrong{Zarr}, preserving multidimensional structure and supporting parallel access:

\begin{sphinxVerbatim}[commandchars=\\\{\}]
\PYG{k+kn}{from}\PYG{+w}{ }\PYG{n+nn}{hsa}\PYG{n+nn}{.}\PYG{n+nn}{compute}\PYG{+w}{ }\PYG{k+kn}{import} \PYG{n}{write\PYGZus{}raster\PYGZus{}stack\PYGZus{}zarr}

\PYG{n}{write\PYGZus{}raster\PYGZus{}stack\PYGZus{}zarr}\PYG{p}{(}
    \PYG{n}{env}\PYG{p}{,}
    \PYG{l+s+s2}{\PYGZdq{}}\PYG{l+s+s2}{environment.zarr}\PYG{l+s+s2}{\PYGZdq{}}\PYG{p}{,}
    \PYG{n}{target\PYGZus{}chunk\PYGZus{}mb}\PYG{o}{=}\PYG{l+m+mi}{128}\PYG{p}{,}
\PYG{p}{)}
\end{sphinxVerbatim}

\sphinxAtStartPar
Large sampled point tables may be stored as \sphinxstylestrong{Parquet}:

\begin{sphinxVerbatim}[commandchars=\\\{\}]
\PYG{k+kn}{from}\PYG{+w}{ }\PYG{n+nn}{hsa}\PYG{n+nn}{.}\PYG{n+nn}{compute}\PYG{+w}{ }\PYG{k+kn}{import} \PYG{n}{write\PYGZus{}table\PYGZus{}parquet}

\PYG{n}{write\PYGZus{}table\PYGZus{}parquet}\PYG{p}{(}\PYG{n}{samples}\PYG{p}{,} \PYG{l+s+s2}{\PYGZdq{}}\PYG{l+s+s2}{samples.parquet}\PYG{l+s+s2}{\PYGZdq{}}\PYG{p}{)}
\end{sphinxVerbatim}

\sphinxAtStartPar
These intermediate products separate expensive environmental processing from statistical fitting. An analysis can therefore be repeated with a different model specification without downloading, transforming and sampling the complete environmental dataset again.


\section{Chunking and Memory Management}
\label{\detokenize{implementation:chunking-and-memory-management}}
\sphinxAtStartPar
High\sphinxhyphen{}resolution analyses commonly exceed available memory long before they exceed available computing power. \sphinxstylestrong{hrHSA} therefore partitions raster and point operations into manageable blocks.

\sphinxAtStartPar
For a raster containing \(B\) bands, spatial chunk dimensions \(n_x\) and \(n_y\), and elements requiring \(d\) bytes, the approximate memory occupied by one chunk is
\begin{equation*}
\begin{split}
M_{\mathrm{chunk}}
=
B n_x n_y d.
\end{split}
\end{equation*}
\sphinxAtStartPar
The package provides heuristics for selecting approximately square spatial chunks under a specified memory target. Point extraction is similarly divided into batches whose size depends upon the number of environmental bands and their numeric precision.

\sphinxAtStartPar
These heuristics provide conservative starting values rather than universal optima. Suitable chunk sizes depend upon the storage system, network throughput, worker memory and downstream operation. Computationally demanding analyses should therefore be benchmarked under realistic conditions.


\section{Distributed and HPC Execution}
\label{\detokenize{implementation:distributed-and-hpc-execution}}
\sphinxAtStartPar
Distributed execution is provided through \sphinxstylestrong{Dask}. The statistical and geospatial functions remain independent of the scheduler, allowing the same analysis to operate under different computational backends.

\sphinxAtStartPar
A local client may be created for notebook or workstation analyses:

\begin{sphinxVerbatim}[commandchars=\\\{\}]
\PYG{k+kn}{from}\PYG{+w}{ }\PYG{n+nn}{hsa}\PYG{n+nn}{.}\PYG{n+nn}{compute}\PYG{+w}{ }\PYG{k+kn}{import} \PYG{n}{make\PYGZus{}local\PYGZus{}dask\PYGZus{}client}

\PYG{n}{client} \PYG{o}{=} \PYG{n}{make\PYGZus{}local\PYGZus{}dask\PYGZus{}client}\PYG{p}{(}
    \PYG{n}{n\PYGZus{}workers}\PYG{o}{=}\PYG{l+m+mi}{8}\PYG{p}{,}
    \PYG{n}{threads\PYGZus{}per\PYGZus{}worker}\PYG{o}{=}\PYG{l+m+mi}{1}\PYG{p}{,}
    \PYG{n}{local\PYGZus{}directory}\PYG{o}{=}\PYG{l+s+s2}{\PYGZdq{}}\PYG{l+s+s2}{/tmp/hsa\PYGZhy{}dask}\PYG{l+s+s2}{\PYGZdq{}}\PYG{p}{,}
\PYG{p}{)}
\end{sphinxVerbatim}

\sphinxAtStartPar
On a SLURM\sphinxhyphen{}managed cluster, the computational backend may be replaced without altering the scientific workflow:

\begin{sphinxVerbatim}[commandchars=\\\{\}]
\PYG{k+kn}{from}\PYG{+w}{ }\PYG{n+nn}{dask}\PYG{n+nn}{.}\PYG{n+nn}{distributed}\PYG{+w}{ }\PYG{k+kn}{import} \PYG{n}{Client}
\PYG{k+kn}{from}\PYG{+w}{ }\PYG{n+nn}{hsa}\PYG{n+nn}{.}\PYG{n+nn}{compute}\PYG{+w}{ }\PYG{k+kn}{import} \PYG{n}{make\PYGZus{}slurm\PYGZus{}cluster}

\PYG{n}{cluster} \PYG{o}{=} \PYG{n}{make\PYGZus{}slurm\PYGZus{}cluster}\PYG{p}{(}
    \PYG{n}{queue}\PYG{o}{=}\PYG{l+s+s2}{\PYGZdq{}}\PYG{l+s+s2}{general}\PYG{l+s+s2}{\PYGZdq{}}\PYG{p}{,}
    \PYG{n}{cores}\PYG{o}{=}\PYG{l+m+mi}{4}\PYG{p}{,}
    \PYG{n}{processes}\PYG{o}{=}\PYG{l+m+mi}{4}\PYG{p}{,}
    \PYG{n}{memory}\PYG{o}{=}\PYG{l+s+s2}{\PYGZdq{}}\PYG{l+s+s2}{32GB}\PYG{l+s+s2}{\PYGZdq{}}\PYG{p}{,}
    \PYG{n}{walltime}\PYG{o}{=}\PYG{l+s+s2}{\PYGZdq{}}\PYG{l+s+s2}{04:00:00}\PYG{l+s+s2}{\PYGZdq{}}\PYG{p}{,}
    \PYG{n}{scale\PYGZus{}jobs}\PYG{o}{=}\PYG{l+m+mi}{10}\PYG{p}{,}
\PYG{p}{)}

\PYG{n}{client} \PYG{o}{=} \PYG{n}{Client}\PYG{p}{(}\PYG{n}{cluster}\PYG{p}{)}
\end{sphinxVerbatim}

\sphinxAtStartPar
Worker\sphinxhyphen{}memory thresholds are configured explicitly so that intermediate data may spill to local storage before workers exhaust their available memory. Large raster stacks can be persisted across workers when they are accessed repeatedly, and temporary data can be directed to node\sphinxhyphen{}local scratch storage.

\sphinxAtStartPar
External services require explicit worker initialisation. Google Earth Engine, for example, maintains process\sphinxhyphen{}local state and must be initialised separately on each distributed worker. \sphinxstylestrong{hrHSA} exposes this operation directly rather than relying upon hidden state inherited from the client process.


\section{Statistical Validation}
\label{\detokenize{implementation:statistical-validation}}
\sphinxAtStartPar
Validation is part of the implementation rather than an optional final step. The package supports temporally blocked cross\sphinxhyphen{}validation and leave\sphinxhyphen{}one\sphinxhyphen{}individual\sphinxhyphen{}out validation, together with Boyce\sphinxhyphen{}index and calibration diagnostics.

\sphinxAtStartPar
Temporal blocks reduce leakage between neighbouring relocations, while leave\sphinxhyphen{}one\sphinxhyphen{}individual\sphinxhyphen{}out validation tests whether a model transfers to animals not represented during fitting. In the latter case, training availability is generated separately within each training animal’s domain, and the held\sphinxhyphen{}out animal is evaluated relative to its own availability.

\sphinxAtStartPar
Validation returns fold\sphinxhyphen{}specific coefficients, predictions, surfaces, Boyce diagnostics, calibration tables and sample sizes. Retaining these outputs permits individual failures and unstable coefficients to be examined instead of concealing them within a single pooled score.

\sphinxAtStartPar
Movement modules provide trajectory preparation, step\sphinxhyphen{}length and turning\sphinxhyphen{}angle calculation, and estimation of empirical movement distributions by individual. These components supply the movement information required by subsequent SSF and iSSF workflows.


\section{Bayesian and Simulation Extensions}
\label{\detokenize{implementation:bayesian-and-simulation-extensions}}
\sphinxAtStartPar
The package architecture is intended to support hierarchical Bayesian inference and large stochastic simulation ensembles. Feature construction, environmental sampling and validation are kept independent of the fitting backend so that likelihood\sphinxhyphen{}based estimation may be supplemented by PyMC\sphinxhyphen{} or JAX\sphinxhyphen{}based models without rebuilding the complete geospatial workflow.

\sphinxAtStartPar
Gradient\sphinxhyphen{}based algorithms such as Hamiltonian Monte Carlo and the No\sphinxhyphen{}U\sphinxhyphen{}Turn Sampler are particularly suitable for the high\sphinxhyphen{}dimensional and correlated parameter spaces produced by hierarchical habitat\sphinxhyphen{}selection models (Hoffman \& Gelman, 2014; Betancourt, 2017). Independent chains can be distributed across processes or nodes, while automatic differentiation and just\sphinxhyphen{}in\sphinxhyphen{}time compilation may accelerate repeated likelihood evaluation.

\sphinxAtStartPar
Forward simulations provide a second level of parallelism. Independent trajectories or posterior simulations may be distributed among workers, while the candidate steps within each trajectory can be evaluated through vectorised array operations.

\sphinxAtStartPar
The present package provides the geospatial, movement and distributed\sphinxhyphen{}computing foundations for these extensions. Fully integrated Bayesian fitting and habitat\sphinxhyphen{}biased forward simulation remain developing components and are kept distinct from the stable RSF implementation.


\section{Documentation and Quality Assurance}
\label{\detokenize{implementation:documentation-and-quality-assurance}}
\sphinxAtStartPar
The documentation is built with \sphinxstylestrong{Sphinx} and \sphinxstylestrong{MyST Markdown}, providing mathematical notation, automatically generated API references and versioned narrative documentation. Package configuration is maintained in \sphinxcode{\sphinxupquote{pyproject.toml}}, with explicit core, optional and development dependencies.

\sphinxAtStartPar
The current test suite verifies the public import surface and provides a foundation for continued quality assurance. Release hardening should additionally include:
\begin{itemize}
\item {} 
\sphinxAtStartPar
unit tests for spatial and statistical functions;

\item {} 
\sphinxAtStartPar
regression tests for stable numerical outputs;

\item {} 
\sphinxAtStartPar
synthetic parameter\sphinxhyphen{}recovery experiments;

\item {} 
\sphinxAtStartPar
end\sphinxhyphen{}to\sphinxhyphen{}end workflow tests;

\item {} 
\sphinxAtStartPar
and benchmarks of runtime, memory use and scaling efficiency.

\end{itemize}

\sphinxAtStartPar
Scientific verification must accompany software verification. A function may execute correctly while implementing an inappropriate availability domain, temporal scale or ecological model.


\section{Reproducibility}
\label{\detokenize{implementation:reproducibility}}
\sphinxAtStartPar
A reproducible \sphinxstylestrong{hrHSA} analysis should preserve:
\begin{itemize}
\item {} 
\sphinxAtStartPar
package and dependency versions;

\item {} 
\sphinxAtStartPar
coordinate reference systems and raster grains;

\item {} 
\sphinxAtStartPar
the definition of availability;

\item {} 
\sphinxAtStartPar
random seeds and sampling ratios;

\item {} 
\sphinxAtStartPar
the feature specification and preprocessing metadata;

\item {} 
\sphinxAtStartPar
fitted models and diagnostic outputs;

\item {} 
\sphinxAtStartPar
cross\sphinxhyphen{}validation partitions;

\item {} 
\sphinxAtStartPar
and the environmental products, or immutable references, from which predictors were derived.

\end{itemize}

\sphinxAtStartPar
Reproducibility is therefore treated as part of the analytical design rather than as an administrative addition made after the analysis is complete.


\section{Summary}
\label{\detokenize{implementation:summary}}
\sphinxAtStartPar
\sphinxstylestrong{hrHSA} combines an accessible Python interface with a modular architecture for high\sphinxhyphen{}resolution telemetry and environmental data. Declarative model specifications, labelled geospatial arrays, memory\sphinxhyphen{}aware batching, chunked storage, structured validation and interchangeable local or SLURM\sphinxhyphen{}backed execution allow the same scientific workflow to operate across different computational scales.

\sphinxAtStartPar
The purpose of this architecture is not parallelism for its own sake. It is to ensure that increasing spatial resolution, study extent and model complexity does not require a corresponding loss of transparency, reproducibility or ecological rigour.


\section{Principal References}
\label{\detokenize{implementation:principal-references}}
\sphinxAtStartPar
Betancourt, M. (2017). A conceptual introduction to Hamiltonian Monte Carlo. \sphinxstyleemphasis{arXiv}, 1701.02434.

\sphinxAtStartPar
Hoffman, M. D. \& Gelman, A. (2014). The No\sphinxhyphen{}U\sphinxhyphen{}Turn Sampler: adaptively setting path lengths in Hamiltonian Monte Carlo. \sphinxstyleemphasis{Journal of Machine Learning Research}, \sphinxstylestrong{15}, 1593\textendash{}1623.

\sphinxAtStartPar
Hoyer, S. \& Hamman, J. J. (2017). xarray: N\sphinxhyphen{}D labelled arrays and datasets in Python. \sphinxstyleemphasis{Journal of Open Research Software}, \sphinxstylestrong{5}, 10.

\sphinxAtStartPar
Rocklin, M. (2015). Dask: parallel computation with blocked algorithms and task scheduling. \sphinxstyleemphasis{Proceedings of the 14th Python in Science Conference}, 130\textendash{}136.

\sphinxAtStartPar
Sandve, G. K., Nekrutenko, A., Taylor, J. \& Hovig, E. (2013). Ten simple rules for reproducible computational research. \sphinxstyleemphasis{PLoS Computational Biology}, \sphinxstylestrong{9}, e1003285.

\sphinxAtStartPar
Wilkinson, M. D. et al. (2016). The FAIR Guiding Principles for scientific data management and stewardship. \sphinxstyleemphasis{Scientific Data}, \sphinxstylestrong{3}, 160018.

\sphinxAtStartPar
Wilson, G. et al. (2017). Good enough practices in scientific computing. \sphinxstyleemphasis{PLoS Computational Biology}, \sphinxstylestrong{13}, e1005510.

\sphinxstepscope


\chapter{Case Studies}
\label{\detokenize{case-studies:case-studies}}\label{\detokenize{case-studies::doc}}

\section{Static RSF (Pangolin)}
\label{\detokenize{case-studies:static-rsf-pangolin}}

\section{Seasonally\sphinxhyphen{}Variant RSF (Lion)}
\label{\detokenize{case-studies:seasonally-variant-rsf-lion}}

\section{Dynamic SSF (Vultures)}
\label{\detokenize{case-studies:dynamic-ssf-vultures}}
\sphinxstepscope


\chapter{API reference}
\label{\detokenize{api:api-reference}}\label{\detokenize{api::doc}}


\renewcommand{\indexname}{Index}
\printindex
\end{document}